% Chapter 5. All quoted numbers come from chapters/ch05-numbers.tex, written by code/ch05/run_all.py.
% Generated by code/ch05/run_all.py. Do not edit by hand.
\newcommand{\NdOne}{0.6368}
\newcommand{\NdTwo}{0.5596}
\newcommand{\acRetAvg}{0.13}
\newcommand{\acRetNoon}{0.00}
\newcommand{\adjCall}{8.4213}
\newcommand{\atmHesLong}{18.11}
\newcommand{\atmHesMid}{18.38}
\newcommand{\atmHesShort}{19.50}
\newcommand{\atmMerLong}{18.95}
\newcommand{\atmMerMid}{18.28}
\newcommand{\atmMerShort}{17.04}
\newcommand{\blackCall}{5.9177}
\newcommand{\blackPut}{7.8781}
\newcommand{\bsApprox}{19.967}
\newcommand{\bsApproxSpot}{26.2}
\newcommand{\bsApproxStrike}{105.13}
\newcommand{\callBase}{10.4506}
\newcommand{\callDeep}{54.9701}
\newcommand{\callDeepLower}{54.8771}
\newcommand{\cnLambda}{222}
\newcommand{\cnSpace}{800}
\newcommand{\cnSteps}{10}
\newcommand{\costDaily}{1.077}
\newcommand{\costFiftyTwo}{0.555}
\newcommand{\costInit}{0.127}
\newcommand{\costRate}{0.2}
\newcommand{\costRebalDaily}{0.950}
\newcommand{\costRebalFiftyTwo}{0.428}
\newcommand{\costRebalThirteen}{0.205}
\newcommand{\costSlope}{0.52}
\newcommand{\costThirteen}{0.332}
\newcommand{\dOne}{0.3500}
\newcommand{\dTwo}{0.1500}
\newcommand{\deltaC}{0.6368}
\newcommand{\deltaLR}{0.6370}
\newcommand{\deltaLRSE}{0.0046}
\newcommand{\deltaP}{\ensuremath{-}0.3632}
\newcommand{\deltaPW}{0.6375}
\newcommand{\deltaPWSE}{0.0018}
\newcommand{\deltaPaths}{100,000}
\newcommand{\deltaSERatio}{2.53}
\newcommand{\deltaVarRatio}{6.4}
\newcommand{\densAtmVol}{18.38}
\newcommand{\densDisc}{0.9753}
\newcommand{\densFwd}{102.532}
\newcommand{\densInt}{1.0000}
\newcommand{\densIntRaw}{0.9753}
\newcommand{\densMean}{102.531}
\newcommand{\densTol}{0.0001}
\newcommand{\divPV}{3.9015}
\newcommand{\divYieldCall}{8.6525}
\newcommand{\divYieldDelta}{0.5621}
\newcommand{\divYieldFwd}{102.0201}
\newcommand{\divYieldPut}{6.7309}
\newcommand{\dropAdjGap}{\ensuremath{-}0.147}
\newcommand{\dropAdjGapHigh}{\ensuremath{-}0.086}
\newcommand{\dropAdjGapLow}{\ensuremath{-}0.020}
\newcommand{\dropCall}{8.2747}
\newcommand{\dropEscGap}{0.161}
\newcommand{\dropGapMax}{0.17}
\newcommand{\dropPairs}{500,000}
\newcommand{\dropPaths}{1,000,000}
\newcommand{\dropSE}{0.0103}
\newcommand{\escCall}{8.1136}
\newcommand{\explAmpC}{1.07}
\newcommand{\explCourC}{0.5182}
\newcommand{\explErrA}{\ensuremath{1.5\times 10^{-4}}}
\newcommand{\explErrB}{0.0042}
\newcommand{\explErrC}{\ensuremath{1.2\times 10^{6}}}
\newcommand{\explStepsA}{309}
\newcommand{\explStepsB}{278}
\newcommand{\explStepsC}{268}
\newcommand{\fdCNValue}{10.4498}
\newcommand{\fdErrCN}{\ensuremath{8.1\times 10^{-4}}}
\newcommand{\fdErrCNFine}{\ensuremath{8.1\times 10^{-4}}}
\newcommand{\fdErrCNR}{\ensuremath{8.3\times 10^{-4}}}
\newcommand{\fdErrExpl}{\ensuremath{5.6\times 10^{-5}}}
\newcommand{\fdErrImpl}{0.0061}
\newcommand{\fdOrderCN}{2.00}
\newcommand{\fdOrderCNR}{2.00}
\newcommand{\fdOrderExpl}{2.01}
\newcommand{\fdOrderImpl}{1.11}
\newcommand{\fdStepRatio}{6.9}
\newcommand{\fdStepsCN}{200}
\newcommand{\fdStepsCNFine}{1,600}
\newcommand{\fdStepsExpl}{1,389}
\newcommand{\fxBest}{135}
\newcommand{\fxBestMonth}{December 2008}
\newcommand{\fxCorrGT}{0.973}
\newcommand{\fxCum}{281}
\newcommand{\fxDays}{4,853}
\newcommand{\fxEnd}{29 September 2026}
\newcommand{\fxFirst}{July 2007}
\newcommand{\fxLast}{August 2026}
\newcommand{\fxLastClose}{29 September 2026}
\newcommand{\fxN}{230}
\newcommand{\fxNAvg}{116}
\newcommand{\fxNNoon}{114}
\newcommand{\fxPnlMean}{1.2}
\newcommand{\fxPnlMeanAvg}{1.4}
\newcommand{\fxPnlMeanExWorst}{3.2}
\newcommand{\fxPnlMeanNoon}{1.0}
\newcommand{\fxPnlMedian}{1.6}
\newcommand{\fxPnlSE}{1.9}
\newcommand{\fxPnlSd}{28.8}
\newcommand{\fxPremMean}{88.2}
\newcommand{\fxRsq}{0.62}
\newcommand{\fxShAbove}{57}
\newcommand{\fxShMean}{7.67}
\newcommand{\fxSlope}{0.72}
\newcommand{\fxSlopeSE}{0.04}
\newcommand{\fxSrMean}{7.53}
\newcommand{\fxStart}{May 2007}
\newcommand{\fxTopFiveShare}{20}
\newcommand{\fxWin}{54}
\newcommand{\fxWindow}{63}
\newcommand{\fxWorst}{\ensuremath{-}101}
\newcommand{\fxWorstLoss}{101}
\newcommand{\fxWorstMonth}{October 2007}
\newcommand{\fxWorstSh}{10.0}
\newcommand{\fxWorstSr}{16.6}
\newcommand{\fxYears}{19}
\newcommand{\gammaB}{0.01876}
\newcommand{\gammaErrCN}{5.0940}
\newcommand{\gammaErrCNR}{0.0001}
\newcommand{\gammaErrImpl}{0.0016}
\newcommand{\gammaPeakExact}{0.0408}
\newcommand{\gammaPeakS}{89.58}
\newcommand{\gkCall}{9,275}
\newcommand{\gkDelta}{0.3052}
\newcommand{\gkFwd}{1.3748}
\newcommand{\gkHedge}{305,206}
\newcommand{\gkPut}{34,284}
\newcommand{\gtApprox}{\ensuremath{-}0.057}
\newcommand{\gtStepCorr}{0.9986}
\newcommand{\gtTotal}{\ensuremath{-}0.093}
\newcommand{\hedgeCorrGTDaily}{0.990}
\newcommand{\hedgeCorrGTFiftyTwo}{0.977}
\newcommand{\hedgeMeanDaily}{0.000}
\newcommand{\hedgeMeanFiftyTwo}{\ensuremath{-}0.009}
\newcommand{\hedgeMeanMaxZ}{2.3}
\newcommand{\hedgeMeanQuarter}{\ensuremath{-}0.001}
\newcommand{\hedgeMeanSEDaily}{0.003}
\newcommand{\hedgeMeanSEFiftyTwo}{0.006}
\newcommand{\hedgeMeanSEQuarter}{0.001}
\newcommand{\hedgeMeanSEThirteen}{0.012}
\newcommand{\hedgeMeanThirteen}{\ensuremath{-}0.028}
\newcommand{\hedgeMu}{10}
\newcommand{\hedgeNDaily}{252}
\newcommand{\hedgeNFiftyTwo}{52}
\newcommand{\hedgeNQuarter}{1,008}
\newcommand{\hedgeNThirteen}{13}
\newcommand{\hedgePaths}{20,000}
\newcommand{\hedgeSdDaily}{0.407}
\newcommand{\hedgeSdFiftyTwo}{0.881}
\newcommand{\hedgeSdPctDaily}{3.9}
\newcommand{\hedgeSdPctFiftyTwo}{8.4}
\newcommand{\hedgeSdPctQuarter}{2.0}
\newcommand{\hedgeSdPctThirteen}{16.4}
\newcommand{\hedgeSdQuarter}{0.207}
\newcommand{\hedgeSdThirteen}{1.719}
\newcommand{\hedgeSlope}{\ensuremath{-}0.487}
\newcommand{\hedgeThDaily}{0.411}
\newcommand{\hedgeThFiftyTwo}{0.906}
\newcommand{\hedgeThGapDaily}{1.2}
\newcommand{\hedgeThQuarter}{0.206}
\newcommand{\hedgeThThirteen}{1.812}
\newcommand{\hesFourier}{5.1808}
\newcommand{\hesMC}{5.1887}
\newcommand{\hesMCPaths}{200,000}
\newcommand{\hesMCSE}{0.0148}
\newcommand{\hesMCSteps}{200}
\newcommand{\hesMCZ}{0.5}
\newcommand{\hesSurfMax}{31.7}
\newcommand{\hesSurfMin}{13.9}
\newcommand{\ivBrentMax}{14}
\newcommand{\ivBrentMed}{9}
\newcommand{\ivExBrentCalls}{11}
\newcommand{\ivExBrentIter}{10}
\newcommand{\ivExErrFour}{\ensuremath{1\times 10^{-14}}}
\newcommand{\ivExIter}{5}
\newcommand{\ivExPrice}{3.0441}
\newcommand{\ivExStart}{0.7933}
\newcommand{\ivExStepA}{0.321209}
\newcommand{\ivExStepB}{0.300347}
\newcommand{\ivExStepC}{0.300000}
\newcommand{\ivGridN}{984}
\newcommand{\ivGridWorst}{\ensuremath{6\times 10^{-15}}}
\newcommand{\ivNewtonMax}{6}
\newcommand{\ivNewtonMed}{5}
\newcommand{\lelandExInit}{0.004}
\newcommand{\lelandMarkup}{0.421}
\newcommand{\lelandMean}{\ensuremath{-}0.122}
\newcommand{\lelandPrice}{10.8712}
\newcommand{\lelandSE}{0.006}
\newcommand{\lelandSd}{0.915}
\newcommand{\lelandVol}{21.12}
\newcommand{\mcBig}{1,000,000}
\newcommand{\mcCorrAnti}{\ensuremath{-}0.5010}
\newcommand{\mcCoverAnti}{94.0}
\newcommand{\mcCoverBoth}{94.0}
\newcommand{\mcCoverCV}{97.0}
\newcommand{\mcCoverPlain}{96.0}
\newcommand{\mcEstAnti}{10.4636}
\newcommand{\mcEstBoth}{10.4545}
\newcommand{\mcEstCV}{10.4519}
\newcommand{\mcEstPlain}{10.4524}
\newcommand{\mcExact}{10.4506}
\newcommand{\mcPairs}{500,000}
\newcommand{\mcReps}{200}
\newcommand{\mcRho}{0.9245}
\newcommand{\mcRhoEven}{0.964}
\newcommand{\mcSEAnti}{0.0104}
\newcommand{\mcSEBoth}{0.0028}
\newcommand{\mcSECV}{0.0056}
\newcommand{\mcSEMillionTheory}{0.0147}
\newcommand{\mcSEPlain}{0.0147}
\newcommand{\mcSdY}{14.719}
\newcommand{\mcSlopeAnti}{\ensuremath{-}0.50}
\newcommand{\mcSlopeBoth}{\ensuremath{-}0.49}
\newcommand{\mcSlopeCV}{\ensuremath{-}0.49}
\newcommand{\mcSlopePlain}{\ensuremath{-}0.51}
\newcommand{\merTotalVol}{19.69}
\newcommand{\noDivCall}{10.4506}
\newcommand{\parityRHS}{4.8771}
\newcommand{\putBase}{5.5735}
\newcommand{\rhoC}{53.2325}
\newcommand{\rhoP}{\ensuremath{-}41.8905}
\newcommand{\rtailHes}{3.47}
\newcommand{\rtailHesGrid}{3.46}
\newcommand{\rtailLogn}{7.11}
\newcommand{\sigAdj}{20.81}
\newcommand{\skewHesLong}{2.90}
\newcommand{\skewHesMid}{6.07}
\newcommand{\skewHesShort}{7.60}
\newcommand{\skewMerLong}{0.71}
\newcommand{\skewMerMid}{2.14}
\newcommand{\skewMerShort}{5.54}
\newcommand{\snipFormula}{10.4506}
\newcommand{\snipMC}{10.4441}
\newcommand{\snipOdds}{55}
\newcommand{\snipPairs}{500,000}
\newcommand{\snipProb}{1.8}
\newcommand{\snipSE}{0.0028}
\newcommand{\snipSeed}{1}
\newcommand{\snipZ}{2.36}
\newcommand{\tailHes}{8.34}
\newcommand{\tailHesGrid}{8.36}
\newcommand{\tailLogn}{4.93}
\newcommand{\thetaC}{\ensuremath{-}6.4140}
\newcommand{\thetaCDay}{\ensuremath{-}0.0176}
\newcommand{\thetaCarryPart}{\ensuremath{-}2.6616}
\newcommand{\thetaGammaPart}{\ensuremath{-}3.7524}
\newcommand{\thetaP}{\ensuremath{-}1.6579}
\newcommand{\thetaPDay}{\ensuremath{-}0.0045}
\newcommand{\underGap}{\ensuremath{-}1.8854}
\newcommand{\underGapAbs}{1.8854}
\newcommand{\underLoss}{1.893}
\newcommand{\underMean}{\ensuremath{-}1.893}
\newcommand{\underSE}{0.007}
\newcommand{\vannaB}{\ensuremath{-}0.2814}
\newcommand{\vegaB}{37.5240}
\newcommand{\vegaPoint}{0.3752}
\newcommand{\volGap}{1.8589}
\newcommand{\volgaB}{9.8501}
\newcommand{\vrAntiTheory}{2.00}
\newcommand{\vrBothTheory}{28.66}
\newcommand{\vrCVTheory}{6.88}
\newcommand{\vrMCAnti}{2.00}
\newcommand{\vrMCBoth}{28.65}
\newcommand{\vrMCCV}{6.88}
\newcommand{\vrMCProduct}{13.8}
\newcommand{\wrongImpLoss}{0.00}
\newcommand{\wrongImpMean}{1.850}
\newcommand{\wrongImpMeanDrift}{1.767}
\newcommand{\wrongImpSE}{0.005}
\newcommand{\wrongImpSEDrift}{0.005}
\newcommand{\wrongImpSd}{0.696}
\newcommand{\wrongImpSdDrift}{0.716}
\newcommand{\wrongRealLoss}{0.02}
\newcommand{\wrongRealMean}{1.855}
\newcommand{\wrongRealSd}{0.310}
\newcommand{\xvEightFirst}{26.5414}
\newcommand{\xvEightIter}{3}
\newcommand{\xvEightSig}{26.5414}
\newcommand{\xvEightStart}{26.48}
\newcommand{\xvFiveFwd}{1.3498}
\newcommand{\xvFivePut}{35,388}
\newcommand{\xvFourApprox}{4.7515}
\newcommand{\xvFourCall}{4.7471}
\newcommand{\xvNineCheck}{0.250}
\newcommand{\xvNineN}{683}
\newcommand{\xvNinePerDay}{2.7}
\newcommand{\xvNineSd}{0.906}
\newcommand{\xvOneCall}{2.0300}
\newcommand{\xvOneDOne}{\ensuremath{-}0.3376}
\newcommand{\xvOneDTwo}{\ensuremath{-}0.5144}
\newcommand{\xvOneParity}{\ensuremath{-}3.9109}
\newcommand{\xvOnePut}{5.9409}
\newcommand{\xvSevenDt}{0.00111}
\newcommand{\xvSevenSteps}{900}
\newcommand{\xvTenGap}{1.8854}
\newcommand{\xvThreeS}{51.44}
\newcommand{\xvTwoApprox}{2.4191}
\newcommand{\xvTwoDelta}{0.3678}
\newcommand{\xvTwoDeltaOnly}{2.3978}
\newcommand{\xvTwoExact}{2.4192}
\newcommand{\xvTwoGamma}{0.0426}

\begin{refsection}
\setcounter{chapter}{4}
\chapter{Black-Scholes: Derivation, Greeks and Numerical Checks}
\label{ch:bs}

Most of the options an alternative investor meets are not listed. A convertible arbitrage fund buys
the equity option buried in a convertible bond, a volatility fund sells index variance, a manager's
carried interest behaves like a call on the fund, and a real estate developer holds an option to wait
before building. These claims are priced, hedged and argued about in one language, the language of the
model that \textcite{black1973} and \textcite{merton1973} wrote down. Even people who think the model
is wrong quote prices in its units, as implied volatilities, and measure their risks with its Greeks.

The chapter follows one question from start to finish: what should a dealer charge for a call, and
what happens to the dealer who charges it? Two derivations give the price, one by building a hedge and
one by taking an expectation under a changed measure. Before relying on the resulting formula we
recompute it by simulation and on a grid, and run it backwards to read volatility from prices. The
last part of the chapter turns to the dealer's profit. A dealer who hedges at the wrong volatility earns
an amount that can be written down exactly, and one who hedges only once a day carries an error whose
size can be predicted; we measure both on simulated paths, and the first on \fxYears{} years of US dollar
options written in Canada.

\begin{keybox}[How to read this chapter]
The numbers in the text are computed by eight scripts in \texttt{code/ch05/}
(\texttt{closed\_form}, \texttt{dividends\_black}, \texttt{monte\_carlo}, \texttt{finite\_diff},
\texttt{implied\_vol}, \texttt{hedging\_pnl}, \texttt{case\_fx\_hedge} and \texttt{exercises}), which
rely on NumPy, SciPy, pandas, statsmodels and Matplotlib for the numerical work and the figures.
The tests in \texttt{code/ch05/test\_*.py} check the formulas against finite differences, simulation
and each other. An independent R script recomputes the Greeks, the finite difference scheme, the
smiles, the hedging error and the whole case study, and GNU Octave checks the prices by quadrature,
the schemes with sparse matrices and the stochastic volatility price with its own integrator. Unless
stated otherwise the examples use $S_0=100$, $K=100$, $r=5\%$, $\sigma=20\%$, $T=1$ and no dividends.
Time to expiry is written $\tau=T-t$.
Every simulation draws from NumPy's default generator seeded with the base seed \texttt{505}
and a stream number fixed in each script, so any script run alone gives the numbers printed here; the short listing in \cref{sec5:mc} uses seed \snipSeed{} so that it stands on its
own. Monte Carlo standard errors are computed from the quantity that is actually averaged: the
pair average for antithetic draws, and the payoff less the fitted multiple of the control for
control variates.

Each symbol keeps one meaning. $\Delta$ on its own is the delta and $\Delta_t$, with a time subscript,
the number of shares a strategy holds; $\Delta x$, $\Delta\tau$ and $\Delta t$, written without a
subscript, are steps in log price, in time to expiry and in calendar time. $\theta$ is the market price of
risk, as in Chapter~3, and $\rho$ is the Greek rho; correlations are written $\varrho$. Total
volatility $\sigma\sqrt\tau$ is $\nu$, and $v_t$ is the variance in the stochastic volatility
model. Jump parameters carry a subscript $J$. The dividend or income yield is $q$, the usual letter
in option pricing; Chapter~4 wrote $\delta$ for the same quantity. $\mathbb P$ is the real-world
probability and $Q$ the risk-neutral one, so that $P$ is free for the put price.
\end{keybox}

% =====================================================================================
\section{What the model assumes about the market}
\label{sec5:model}

\begin{definition}\label{def5:market}
The \emph{Black-Scholes market} on $[0,T]$ has two assets. A bank account grows at a constant rate,
$dB_t=rB_t\,dt$. A stock pays dividends continuously at a constant yield $q$ and its price follows
\begin{equation}\label{eq5:gbm}
  dS_t=\mu S_t\,dt+\sigma S_t\,dW_t,\qquad S_0>0,
\end{equation}
where $W$ is a Brownian motion under the real-world probability $\mathbb P$, $\sigma>0$ and $\mu$ are
constants, and information is the filtration $(\mathcal F_t)$ generated by $W$. There are no taxes or
trading costs, any real number of shares can be held, and borrowing and lending are at the same rate.
\end{definition}

As shown in Chapter~3, $S_t=S_0\exp\bigl((\mu-\tfrac12\sigma^2)t+\sigma W_t\bigr)$, so log returns over
any interval are normal and independent of the past. The total expected return of the stock is
$\mu+q$: $\mu$ from the price and $q$ from the dividends.

\begin{definition}\label{def5:selffin}
A \emph{trading strategy} is a predictable process $\Delta_t$, integrable against $S$, the number of shares held, together with an
initial value $\Pi_0$. The strategy, with value $\Pi_t$, is \emph{self-financing} if whatever is not in shares sits in
the bank account and no money is added or withdrawn:
\begin{equation}\label{eq5:selffin}
  d\Pi_t=\Delta_t\,dS_t+q\Delta_tS_t\,dt+r(\Pi_t-\Delta_tS_t)\,dt .
\end{equation}
The strategy is \emph{admissible} if $\int_0^T\Delta_t^2S_t^2\,dt<\infty$ almost surely and $\Pi_t$ is
bounded below. An \emph{arbitrage} is an admissible self-financing strategy with $\Pi_0\le0$,
$\Pi_T\ge0$ and $\mathbb P(\Pi_T>0)>0$.
\end{definition}

The three terms in \cref{eq5:selffin} are the gain on the shares, the dividends they pay, and interest
on the cash, which is negative when the strategy borrows. The lower bound on $\Pi$ rules out doubling
strategies, which in continuous time can manufacture money from nothing if unlimited losses are
allowed along the way.

% =====================================================================================
\section{A price that a hedge enforces}
\label{sec5:hedge}

The first derivation looks for a function $V(t,S)$ such that holding $V_S$ shares, and keeping the rest
in the bank, reproduces $V$ at every moment. It is the continuous-time version of the one-step
replication on the binomial tree of Chapter~4.

\begin{theorem}[The Black-Scholes equation]\label{thm5:pde}
Let $g$ be a payoff with $0\le g(s)\le c_0(1+s)$ for some constant $c_0$, and let $V(t,S)$ be continuously
differentiable once in $t$ and twice in $S$ on $[0,T)\times(0,\infty)$, continuous on
$[0,T]\times(0,\infty)$, bounded below, with $V_S$ bounded and
\begin{equation}\label{eq5:pde}
  V_t+\tfrac12\sigma^2S^2V_{SS}+(r-q)SV_S-rV=0,\qquad V(T,S)=g(S).
\end{equation}
Then the strategy $\Delta_t=V_S(t,S_t)$ with $\Pi_0=V(0,S_0)$ is self-financing and admissible and
satisfies $\Pi_t=V(t,S_t)$ for all $t$, in particular $\Pi_T=g(S_T)$. If the claim paying $g(S_T)$
at $T$ also trades, then in a market with no arbitrage its price is at most $V(0,S_0)$, and if $V$ is
bounded above as well, its price is exactly $V(0,S_0)$.
\end{theorem}

\begin{proof}
Let $\Pi$ satisfy \cref{eq5:selffin} with $\Delta_t=V_S(t,S_t)$. By It\^o's lemma (Chapter~3)
applied to $V(t,S_t)$, with all derivatives at $(t,S_t)$,
\begin{equation*}
  dV=V_t\,dt+V_S\,dS_t+\tfrac12\sigma^2S_t^2V_{SS}\,dt .
\end{equation*}
Subtracting this from \cref{eq5:selffin}, the $dS_t$ terms cancel because $\Delta_t=V_S$, and
\begin{equation*}
  d(\Pi_t-V)=\bigl(qS_tV_S+r\Pi_t-rS_tV_S-V_t-\tfrac12\sigma^2S_t^2V_{SS}\bigr)dt
  =r(\Pi_t-V)\,dt,
\end{equation*}
where the second step uses \cref{eq5:pde} in the form $V_t+\tfrac12\sigma^2S^2V_{SS}=rV-(r-q)SV_S$.
So $e_t=\Pi_t-V(t,S_t)$ solves $de_t=re_t\,dt$ on $[0,T)$ with $e_0=0$, which forces $e_t=0$, and
continuity of $\Pi$ and $V$ carries this to $t=T$. Since $V_S$ is bounded and $S$ has continuous
paths, $\int_0^T\Delta_t^2S_t^2\,dt<\infty$; since $V$ is bounded below, so is $\Pi_t=V(t,S_t)$, and
the strategy is admissible. If the claim traded at a price
$\bar V>V(0,S_0)$, selling it and running the strategy would start with $\bar V-V(0,S_0)>0$ in the bank and end
with exactly enough to pay $g(S_T)$, an arbitrage. If $\bar V<V(0,S_0)$, buy the claim and run the
opposite strategy, holding $-V_S$ shares; its share-and-cash value is $-V(t,S_t)$, which is bounded
below exactly when $V$ is bounded above. In that case the opposite strategy is admissible and the
trade is an arbitrage.
\end{proof}

The hypotheses on $V$ have to be checked claim by claim. For a call or a put, $V$ is positive by
\cref{thm5:bs} and $|V_S|\le e^{-q\tau}$ by \cref{prop5:greeks}, hence $|V_S|\le1$ when $q\ge0$; more generally, for any payoff $g\ge0$, the
Feynman-Kac step after \cref{thm5:rn} writes $V(t,S)=e^{-r\tau}\E_Q[g(S_T)\mid S_t=S]\ge0$, so
nonnegativity supplies the lower bound, while the smoothness, growth and delta-bound hypotheses must
still be checked separately. A European put is also bounded above, by $Ke^{-r\tau}$ and therefore by $K$ when $r\ge0$, so its
price is pinned from both sides. A call is not bounded above, and the theorem gives only that its price cannot exceed
$V(0,S_0)$. Equality for the call then follows from put-call parity, \cref{eq5:parity}, whose proof
uses only static trades in the stock and the bank account: a call priced below the formula, combined
with the put and the static position, would give a put priced below its formula.

Neither the expected return $\mu$ nor any attitude to risk appears in \cref{eq5:pde}. Two
investors who disagree about how fast the stock will grow still agree on the price, because the hedge
cancels the exposure to $dW$ and with it everything that depends on $\mu$. This is the same effect
as in the binomial model, where the probability of an up move did not enter the price. The next result
shows that the equation is forced: any hedge of this form must be the delta hedge.

\begin{proposition}[The hedge is unique]\label{prop5:necessary}
Suppose a self-financing strategy holds $\Delta_t=h(t,S_t)$ shares, with $h$ continuous on
$[0,T)\times(0,\infty)$, and its value is $\Pi_t=V(t,S_t)$ for a function $V$ with the smoothness of
\cref{thm5:pde}. Then $h=V_S$ and $V$
satisfies the partial differential equation in \cref{eq5:pde} on $(0,T)\times(0,\infty)$.
\end{proposition}

\begin{proof}
By It\^o's lemma and \cref{eq5:gbm},
$dV=\bigl(V_t+\mu SV_S+\tfrac12\sigma^2S^2V_{SS}\bigr)dt+\sigma SV_S\,dW_t$, while by
\cref{eq5:selffin}, $d\Pi=\bigl((\mu+q)h S+r(V-hS)\bigr)dt+\sigma hS\,dW_t$. The two are the same It\^o
process, and an It\^o process determines its $dt$ and $dW$ coefficients uniquely up to sets of measure
zero (a continuous martingale of finite variation is constant). Matching the $dW$ terms gives
$\sigma S_th=\sigma S_tV_S$, so $h=V_S$ because $\sigma S_t>0$. Matching the $dt$ terms then gives
$V_t+\mu SV_S+\tfrac12\sigma^2S^2V_{SS}=(\mu+q)SV_S+rV-rSV_S$, in which $\mu SV_S$ cancels and what is
left is \cref{eq5:pde}. The identities hold at $(t,S_t)$ for almost every $t$ and almost every path;
since $S_t$ has a positive density on $(0,\infty)$ for each $t>0$ and both sides are continuous, they
hold everywhere.
\end{proof}

% =====================================================================================
\section{The same price as an expectation}
\label{sec5:rn}

The second derivation changes the probability measure instead of solving an equation. Under a measure
in which the discounted gains process of every traded asset, its discounted price plus the discounted
dividends paid so far, is a martingale, a replicable claim is priced by a discounted conditional
expectation.

\begin{theorem}[Risk-neutral pricing]\label{thm5:rn}
Let $\theta=(\mu+q-r)/\sigma$ and define $Q$ on $\mathcal F_T$ by $dQ/d\mathbb P=\exp(-\theta W_T-\tfrac12\theta^2T)$.
Then
\begin{enumerate}[label=(\roman*)]
  \item $\widetilde W_t=W_t+\theta t$ is a $Q$-Brownian motion and
  $dS_t=(r-q)S_t\,dt+\sigma S_t\,d\widetilde W_t$;
  \item the discounted value of any self-financing strategy satisfies
  \begin{equation*}
    d(e^{-rt}\Pi_t)=e^{-rt}\Delta_t\sigma S_t\,d\widetilde W_t ,
  \end{equation*}
  so it is a $Q$-local martingale;
  \item for a claim $H$ that is $\mathcal F_T$-measurable with $\E_Q[H^2]<\infty$ there is a
  self-financing strategy with $\Pi_T=H$ and
  \begin{equation}\label{eq5:rn}
    \Pi_t=e^{-r(T-t)}\E_Q\bigl[H\bigm|\mathcal F_t\bigr],\qquad 0\le t\le T ,
  \end{equation}
  whose holding satisfies $\int_0^T\Delta_t^2S_t^2\,dt<\infty$ almost surely. If in addition
  $H\ge0$, then $\Pi_t\ge0$ and the strategy is admissible.
\end{enumerate}
\end{theorem}

\begin{proof}
(i) is Girsanov's theorem from Chapter~3 with this $\theta$, and substituting
$dW_t=d\widetilde W_t-\theta\,dt$ in \cref{eq5:gbm} gives $\mu-\sigma\theta=r-q$ as the drift.
(ii) By the product rule and \cref{eq5:selffin},
\begin{equation*}
  d(e^{-rt}\Pi_t)=e^{-rt}\bigl(d\Pi_t-r\Pi_t\,dt\bigr)=e^{-rt}\Delta_t\bigl(dS_t+qS_t\,dt-rS_t\,dt\bigr)
  =e^{-rt}\Delta_t\sigma S_t\,d\widetilde W_t ,
\end{equation*}
using (i) in the last step. (iii) Let $\Lambda_t=\E_Q[e^{-rT}H\mid\mathcal F_t]$, a square-integrable
$Q$-martingale. Because $\theta$ is constant, $\widetilde W$ and $W$ generate the same filtration, and
the martingale representation theorem for Brownian filtrations, which we quote
\parencite[Theorem 5.3.1]{shreve2004}, gives a $\zeta$, which can be taken predictable, with $\E_Q\int_0^T\zeta_t^2\,dt<\infty$
and $\Lambda_t=\Lambda_0+\int_0^t\zeta_u\,d\widetilde W_u$. Take $\Pi_0=\Lambda_0$ and $\Delta_t=e^{rt}\zeta_t/(\sigma S_t)$,
which is predictable, as \cref{def5:selffin} requires, because $\zeta$ is and $S$ is continuous.
By (ii), $e^{-rt}\Pi_t=\Lambda_0+\int_0^t\zeta_u\,d\widetilde W_u=\Lambda_t$, which is \cref{eq5:rn}, and at
$t=T$ it gives $\Pi_T=H$. The holding has $\Delta_t^2S_t^2=e^{2rt}\zeta_t^2/\sigma^2$, whose integral is
finite almost surely because its $Q$-expectation is, and $Q$ and $\mathbb P$ have the same null sets.
If $H\ge0$, \cref{eq5:rn} is a conditional expectation of a nonnegative variable, so $\Pi_t\ge0$.
\end{proof}

For a claim $H=g(S_T)$, the Markov property of $S$ under $Q$ turns \cref{eq5:rn} into a function of
$(t,S_t)$, namely $e^{-r\tau}\E_Q[g(S_T)\mid S_t=s]$ evaluated at $s=S_t$. The Feynman-Kac theorem of
Chapter~3, applied to the $Q$-dynamics in (i), says that a solution $V$ of \cref{eq5:pde} with the
properties of \cref{thm5:pde} has exactly this representation.\footnote{Two hypotheses of that theorem
are weakened, as its remark allows: the state space is $(0,\infty)$, and $V$ is smooth only before $T$
and continuous up to it, because the payoff has a kink. The integrability condition holds because
$V_S$ is bounded and $\E_Q\int_0^TS_t^2\,dt<\infty$, and the limit at $T$ uses the growth bound
$g(s)\le c_0(1+s)$.} So the two derivations agree. The hedge shows why the price cannot be anything
else; the expectation shows how to compute it, and it also works for path-dependent claims that have
no equation in $(t,S)$ alone.

Under $Q$ the drift $\mu$ is replaced by $r-q$, and $\theta$, the market price of risk, is folded
into the probabilities. A stock that a forecaster expects to return 10\% a year has, under $Q$, an
expected total return equal to that of the bank account, dividends included, which is why discounting
at $r$ is correct. $Q$ is a device for pricing and says nothing about the stock's actual expected
return.

% =====================================================================================
\section{Solving for the price of a call and a put}
\label{sec5:formula}

\begin{lemma}\label{lem5:lognormal}
If $U\sim\N(m,\eta^2)$ with $\eta>0$ and $\ell\in\mathbb R$, then $\Pr(U>\ell)=\Phi\bigl((m-\ell)/\eta\bigr)$ and
\begin{equation}\label{eq5:partial}
  \E\bigl[e^U\mathbf 1\{U>\ell\}\bigr]=e^{m+\eta^2/2}\,\Phi\Bigl(\frac{m+\eta^2-\ell}{\eta}\Bigr).
\end{equation}
\end{lemma}

\begin{proof}
The first claim is the definition of $\Phi$ after standardising. For the second, complete the square,
$x-(x-m)^2/(2\eta^2)=m+\eta^2/2-(x-m-\eta^2)^2/(2\eta^2)$. Hence
\begin{equation*}
  \E\bigl[e^U\mathbf 1\{U>\ell\}\bigr]=e^{m+\eta^2/2}\int_\ell^\infty\frac{1}{\eta\sqrt{2\pi}}\,e^{-(x-m-\eta^2)^2/(2\eta^2)}\,dx ,
\end{equation*}
and the integral is the probability that an $\N(m+\eta^2,\eta^2)$ variable exceeds $\ell$.
\end{proof}

\begin{theorem}[Black-Scholes-Merton formula]\label{thm5:bs}
In the market of \cref{def5:market} the prices at time $t$ of European calls and puts with strike $K$
and expiry $T$ are $C(t,S_t)$ and $P(t,S_t)$, where, with $\tau=T-t$,
\begin{equation}\label{eq5:bs}
\begin{aligned}
  C(t,S)&=Se^{-q\tau}\Phi(d_1)-Ke^{-r\tau}\Phi(d_2),\qquad
  P(t,S)=Ke^{-r\tau}\Phi(-d_2)-Se^{-q\tau}\Phi(-d_1),\\
  d_1&=\frac{\ln(S/K)+(r-q+\tfrac12\sigma^2)\tau}{\sigma\sqrt\tau},\qquad d_2=d_1-\sigma\sqrt\tau .
\end{aligned}
\end{equation}
Both functions solve \cref{eq5:pde} with their payoffs as terminal values.
\end{theorem}

\begin{proof}
By \cref{thm5:rn}(i), conditional on $S_t=S$, $U=\ln S_T\sim\N(m,\eta^2)$ under $Q$ with
$m=\ln S+(r-q-\tfrac12\sigma^2)\tau$ and $\eta^2=\sigma^2\tau$. Then
$C=e^{-r\tau}\bigl(\E_Q[e^U\mathbf 1\{U>\ln K\}]-KQ(U>\ln K)\bigr)$. With $\ell=\ln K$,
$(m-\ell)/\eta=d_2$, $(m+\eta^2-\ell)/\eta=d_1$ and $e^{m+\eta^2/2}=Se^{(r-q)\tau}$, so \cref{lem5:lognormal}
gives the call formula. The put follows from put-call parity below, together with
$\Phi(-x)=1-\Phi(x)$. That $C$ solves \cref{eq5:pde} is checked in \cref{prop5:gammatheta}, and the
terminal condition holds because $d_1,d_2\to\pm\infty$ as $\tau\downarrow0$ according as $S>K$ or $S<K$;
at $S=K$ both $\Phi(d_1)$ and $\Phi(d_2)$ tend to $\tfrac12$, so $C\to\tfrac12(S-K)=0$.
\end{proof}

For the base case the formula gives $d_1=\dOne$, $d_2=\dTwo$, $\Phi(d_1)=\NdOne$ and $\Phi(d_2)=\NdTwo$, so
the call is worth \callBase{} and the put \putBase. The put is the European benchmark against which
Chapter~4 measured the early exercise premium, and the same chapter showed that binomial
trees converge to these formulas; here they are reached directly.

\begin{proposition}[Put-call parity and bounds]\label{prop5:parity}
For European options on a stock with a known dividend yield $q$, in any market without arbitrage that
has the stock and a bank account at the rate $r$, and in which dividends can be reinvested in the stock,
\begin{equation}\label{eq5:parity}
  C-P=Se^{-q\tau}-Ke^{-r\tau}.
\end{equation}
In the Black-Scholes market, in addition, $(Se^{-q\tau}-Ke^{-r\tau})^+<C<Se^{-q\tau}$, the price increases
strictly in $\sigma$, it tends to the lower bound as $\sigma\downarrow0$ and to the upper bound as
$\sigma\to\infty$.
\end{proposition}

\begin{proof}
A long call and a short put pay $(S_T-K)^+-(K-S_T)^+=S_T-K$ at $T$. So does holding $e^{-q\tau}$ shares,
with the dividends reinvested in shares, and borrowing $Ke^{-r\tau}$; the law of one price gives
\cref{eq5:parity}. For the bounds, $C=e^{-r\tau}\E_Q[(S_T-K)^+]$ and Jensen's inequality give
$C\ge e^{-r\tau}(\E_QS_T-K)^+=(Se^{-q\tau}-Ke^{-r\tau})^+$, strictly because $(x-K)^+$ is not linear on
the support of $S_T$; and $(S_T-K)^+<S_T$ gives $C<e^{-r\tau}\E_QS_T=Se^{-q\tau}$. Strict monotonicity
in $\sigma$ follows from $\partial C/\partial\sigma>0$ in \cref{prop5:greeks}. As $\sigma\downarrow0$,
$d_1$ and $d_2$ tend to $\pm\infty$ with the sign of $\ln(S/K)+(r-q)\tau$, which gives the lower bound
(if that quantity is zero, $d_1$ and $d_2$ tend to $0$ and
$C\to\tfrac12(Se^{-q\tau}-Ke^{-r\tau})=0$, which is again the lower bound);
as $\sigma\to\infty$, $d_1\to\infty$ and $d_2\to-\infty$, which gives $Se^{-q\tau}$.
\end{proof}

Parity needs no model for how the stock moves, only the stock, the bank account and a known yield. In the base case
$S-Ke^{-rT}=\parityRHS$, which is exactly $\callBase-\putBase$. A deep in-the-money call, $S=150$, is
worth \callDeep, a little above its lower bound of \callDeepLower: almost all of its value is the
forward it contains.

\begin{figure}[t]
  \centering
  \includegraphics[width=\textwidth]{ch05/fig01_price}
  \caption{(a) The call price against the stock price at three times to expiry, with the payoff.
  (b) The delta, which steepens into a step at the strike as expiry approaches. $K=100$, $r=5\%$,
  $\sigma=20\%$.}
  \label{fig5:price}
\end{figure}

% =====================================================================================
\section{How the price moves with its inputs}
\label{sec5:greeks}

The partial derivatives of the price are called Greeks. They measure exposure to the inputs and,
through \cref{thm5:pde}, they are also the hedge. Computing them is short work once one relation
between the two normal densities in the formula is in hand.

\begin{lemma}\label{lem5:phi}
With $\varphi$ the standard normal density and $d_1,d_2$ as in \cref{eq5:bs},
$Se^{-q\tau}\varphi(d_1)=Ke^{-r\tau}\varphi(d_2)$.
\end{lemma}

\begin{proof}
$\varphi(d_2)/\varphi(d_1)=\exp\bigl(\tfrac12(d_1^2-d_2^2)\bigr)=\exp\bigl(\tfrac12(d_1-d_2)(d_1+d_2)\bigr)$.
Here $d_1-d_2=\sigma\sqrt\tau$ and $d_1+d_2=2\bigl(\ln(S/K)+(r-q)\tau\bigr)/(\sigma\sqrt\tau)$, so the
ratio is $(S/K)e^{(r-q)\tau}$, which rearranges to the claim.
\end{proof}

\begin{proposition}[The Greeks]\label{prop5:greeks}
For the call in \cref{eq5:bs}, writing $\Theta=\partial C/\partial t=-\partial C/\partial\tau$,
\begin{align*}
  \Delta&=\frac{\partial C}{\partial S}=e^{-q\tau}\Phi(d_1), &
  \Gamma&=\frac{\partial^2C}{\partial S^2}=\frac{e^{-q\tau}\varphi(d_1)}{S\sigma\sqrt\tau}, &
  \mathcal V&=\frac{\partial C}{\partial\sigma}=Se^{-q\tau}\varphi(d_1)\sqrt\tau,\\
  \rho&=\frac{\partial C}{\partial r}=K\tau e^{-r\tau}\Phi(d_2), &
  \Theta&=-\frac{Se^{-q\tau}\varphi(d_1)\sigma}{2\sqrt\tau}+qSe^{-q\tau}\Phi(d_1)-rKe^{-r\tau}\Phi(d_2). \span\span
\end{align*}
For the put, $\Delta_P=\Delta-e^{-q\tau}$, $\rho_P=\rho-K\tau e^{-r\tau}$,
$\Theta_P=\Theta-qSe^{-q\tau}+rKe^{-r\tau}$, and $\Gamma$ and $\mathcal V$ are the same as for the call.
Finally, $\mathcal V=\sigma\tau S^2\Gamma$.
\end{proposition}

\begin{proof}
Differentiate \cref{eq5:bs} by the chain rule. For any of the inputs $\alpha$,
\begin{equation*}
  \frac{\partial C}{\partial \alpha}=\frac{\partial(Se^{-q\tau})}{\partial \alpha}\Phi(d_1)-\frac{\partial(Ke^{-r\tau})}{\partial \alpha}\Phi(d_2)
  +Se^{-q\tau}\varphi(d_1)\frac{\partial d_1}{\partial \alpha}-Ke^{-r\tau}\varphi(d_2)\frac{\partial d_2}{\partial \alpha},
\end{equation*}
and by \cref{lem5:phi} the last two terms combine to
\begin{equation*}
  Se^{-q\tau}\varphi(d_1)\,\frac{\partial(d_1-d_2)}{\partial \alpha}=Se^{-q\tau}\varphi(d_1)\,\frac{\partial(\sigma\sqrt\tau)}{\partial \alpha}.
\end{equation*}
For $\alpha=S$ and $\alpha=r$ this vanishes,
which leaves $\Delta$ and $\rho$. For $\alpha=\sigma$ the first two terms vanish and the last is
$Se^{-q\tau}\varphi(d_1)\sqrt\tau$. For $\alpha=\tau$ the first two terms are
$-qSe^{-q\tau}\Phi(d_1)+rKe^{-r\tau}\Phi(d_2)$
and the last is $Se^{-q\tau}\varphi(d_1)\sigma/(2\sqrt\tau)$; changing sign gives $\Theta$. Then
$\Gamma=\partial\Delta/\partial S=e^{-q\tau}\varphi(d_1)\,\partial d_1/\partial S$ with
$\partial d_1/\partial S=1/(S\sigma\sqrt\tau)$. The put Greeks follow by differentiating
\cref{eq5:parity}, whose right side does not depend on $\sigma$ and is linear in $S$. Finally
$\sigma\tau S^2\Gamma=\sigma\tau S^2e^{-q\tau}\varphi(d_1)/(S\sigma\sqrt\tau)=\mathcal V$.
\end{proof}

\begin{proposition}[Gamma and theta]\label{prop5:gammatheta}
The call and put prices satisfy $\Theta+\tfrac12\sigma^2S^2\Gamma+(r-q)S\Delta-rV=0$, which is
\cref{eq5:pde}. Equivalently, for a long option that is delta hedged and financed at $r$, the time
decay net of financing is fixed by the gamma: $\Theta-rV+rS\Delta-qS\Delta=-\tfrac12\sigma^2S^2\Gamma$.
\end{proposition}

\begin{proof}
Substitute \cref{prop5:greeks}. The gamma term $\tfrac12\sigma^2S^2\Gamma=Se^{-q\tau}\varphi(d_1)\sigma/(2\sqrt\tau)$
cancels the first term of $\Theta$. The $\Phi(d_1)$ terms, $qSe^{-q\tau}\Phi(d_1)$ from $\Theta$,
$(r-q)Se^{-q\tau}\Phi(d_1)$ from $(r-q)S\Delta$ and $-rSe^{-q\tau}\Phi(d_1)$ from $-rV$, sum to zero. So do the
$\Phi(d_2)$ terms, $-rKe^{-r\tau}\Phi(d_2)$ from $\Theta$ and $+rKe^{-r\tau}\Phi(d_2)$ from $-rV$. For the put use parity, since
$Se^{-q\tau}-Ke^{-r\tau}$ also solves the equation.
\end{proof}

The identity is deterministic: it holds at every point $(t,S)$. What it says about money is that over
a short interval the hedged long option loses $\tfrac12\sigma^2S^2\Gamma\,dt$ to time and earns
$\tfrac12\Gamma(dS)^2$ from the move in the stock. Under the model $(dS)^2=\sigma^2S^2\,dt$, so the
two cancel. With discrete rebalancing the leading-order difference has conditional mean zero under
the model, but the realised hedge error varies from path to path, which is the subject of
\cref{sec5:pnl}.

For the base call, $\Delta=\deltaC$, $\Gamma=\gammaB$, vega is \vegaB, or \vegaPoint{} per percentage
point of volatility, and $\rho=\rhoC$. Theta is \thetaC{} a year, or \thetaCDay{} a calendar day. The
put has $\Delta_P=\deltaP$, $\Theta_P=\thetaP$ and $\rho_P=\rhoP$. \Cref{prop5:gammatheta} splits the
call's theta into $-\tfrac12\sigma^2S^2\Gamma=\thetaGammaPart$, the cost of holding convexity, and
$rV-rS\Delta=\thetaCarryPart$, the interest on the hedge. The second order cross sensitivities are
vanna, $\partial^2C/\partial S\partial\sigma=-e^{-q\tau}\varphi(d_1)d_2/\sigma=\vannaB$, and volga,
$\partial^2C/\partial\sigma^2=\mathcal Vd_1d_2/\sigma=\volgaB$. Each Greek agrees with a central
difference of the price to within the difference's own error, and the tests check this at 200 random
points for calls and puts.

\begin{proposition}[Where gamma peaks]\label{prop5:gammapeak}
For fixed $\tau$, $\Gamma$ as a function of $S$ has a unique maximum at
$S^*=K\exp\bigl(-(r-q+\tfrac32\sigma^2)\tau\bigr)$.
\end{proposition}

\begin{proof}
Up to a constant, $\ln\Gamma=-\tfrac12d_1^2-\ln S$, and $\partial d_1/\partial S=1/(S\sigma\sqrt\tau)$,
so
\begin{equation*}
  \frac{\partial\ln\Gamma}{\partial S}=-\frac1S\Bigl(\frac{d_1}{\sigma\sqrt\tau}+1\Bigr).
\end{equation*}
This is positive when
$d_1<-\sigma\sqrt\tau$ and negative when $d_1>-\sigma\sqrt\tau$, so the unique maximum is where
$d_1=-\sigma\sqrt\tau$, that is $\ln(S/K)=-(r-q+\tfrac32\sigma^2)\tau$.
\end{proof}

For the one year option gamma peaks at $S^*=\gammaPeakS$, below the strike. As expiry approaches $S^*$
moves up to $K$ and the peak becomes a spike, which \cref{fig5:dg} shows together with the delta
surface. \Cref{fig5:vt} shows vega, largest for long dated options near the money, and theta, which
for a near-the-money option about to expire falls steeply: the time value that remains is spent in the
last few days.

\begin{figure}[t]
  \centering
  \includegraphics[width=\textwidth]{ch05/fig02_delta_gamma}
  \caption{Delta and gamma of a call with $K=100$, $r=5\%$ and $\sigma=20\%$ against the stock price
  and the time to expiry.}
  \label{fig5:dg}
\end{figure}

\begin{figure}[t]
  \centering
  \includegraphics[width=\textwidth]{ch05/fig03_vega_theta}
  \caption{Vega and theta of the same call. Vega is greatest for long dated options near the money;
  theta is most negative for short dated options at the money.}
  \label{fig5:vt}
\end{figure}

% =====================================================================================
\section{Indices, futures and the Canadian dollar}
\label{sec5:div}

Everything so far allowed a continuous dividend yield $q$. Under further assumptions, stated with
each case, the same formula prices options on an index, on a futures contract and on a currency. A
stock that pays a few fixed cash dividends needs separate treatment.

\begin{proposition}[Forward form]\label{prop5:yield}
Let $F=Se^{(r-q)\tau}$ be the forward price for delivery at $T$. Then \cref{eq5:bs} can be written
\begin{equation}\label{eq5:fwdform}
  C=e^{-r\tau}\bigl(F\Phi(d_1)-K\Phi(d_2)\bigr),\quad P=e^{-r\tau}\bigl(K\Phi(-d_2)-F\Phi(-d_1)\bigr),\quad
  d_{1,2}=\frac{\ln(F/K)\pm\tfrac12\sigma^2\tau}{\sigma\sqrt\tau}.
\end{equation}
The price depends on $S$, $r$ and $q$ only through $F$ and the discount factor $e^{-r\tau}$.
\end{proposition}

\begin{proof}
$Se^{-q\tau}=e^{-r\tau}F$ and $\ln(S/K)+(r-q)\tau=\ln(F/K)$; substitute into \cref{eq5:bs}.
\end{proof}

An index with a 3\% yield and the base inputs has forward \divYieldFwd, call \divYieldCall, put
\divYieldPut{} and delta \divYieldDelta. The yield moves value from the call to the put, because the
holder of the call does not receive the dividends.

\begin{proposition}[Futures and currencies]\label{prop5:black}
\begin{enumerate}[label=(\roman*)]
  \item Suppose the futures price is a geometric Brownian motion, $dF_t=\mu_FF_t\,dt+\sigma F_t\,dW_t$,
  the contract is marked to market continuously, variation margin can be borrowed or lent without
  limit at the constant rate $r$, and there are no margin costs. Then an option on it that is settled
  at $T$ satisfies \cref{eq5:pde} with $q=r$, and its price is \cref{eq5:fwdform} with $F$ the current
  futures price \parencite{black1976}.
  \item If $S$ is the price in domestic currency of one unit of a foreign currency, $r_d$ is the
  domestic rate and $r_f$ the foreign rate, an option on the currency satisfies \cref{eq5:pde} with
  $r=r_d$ and $q=r_f$ \parencite{garman1983}.
\end{enumerate}
\end{proposition}

\begin{proof}
(i) A futures position costs nothing to open and its gains are paid into the account as they arise,
so a strategy holding $\Delta_t$ contracts and keeping its value $\Pi_t$ in the bank has
$d\Pi_t=\Delta_t\,dF_t+r\Pi_t\,dt$. This is \cref{eq5:selffin} with $S$ replaced by $F$ and $q=r$, since
then $qS\Delta+r(\Pi-\Delta S)=r\Pi$. The proof of \cref{thm5:pde} goes through unchanged, and so does
\cref{thm5:bs}; with $q=r$, the forward of $F$ is $F$ itself. Unlike a stock, a futures contract has
zero value when opened and pays settlement gains; the equivalence is between the two pricing
equations under the funding assumptions stated, and daily rather than continuous settlement adds a
small error that the proof ignores. (ii) A unit of foreign currency held in a
foreign bank account earns $r_f$, paid in units of the currency, exactly as a share earns a dividend
yield. So holding $\Delta$ units gives \cref{eq5:selffin} with $q=r_f$ and $r=r_d$.
\end{proof}

A call on a futures contract at 70, struck at 72, with six months to run, $r=4\%$ and $\sigma=35\%$, is
worth \blackCall, and the put \blackPut. For a currency, take a three month call on one million US
dollars with spot 1.38 Canadian dollars per US dollar, strike 1.40, $r_d=2.5\%$, $r_f=4\%$ and
$\sigma=7\%$. The forward is \gkFwd, below spot because the US rate is higher, and the call costs
\gkCall{} Canadian dollars against \gkPut{} for the put. The spot delta is \gkDelta{} Canadian dollars
of option value per one Canadian dollar rise in the price of a US dollar, per US dollar of notional, so
the hedge for the call on one million US dollars is to hold \gkHedge{} US dollars in a US account,
where they earn $r_f$.

\paragraph{Cash dividends}
A stock that pays a fixed amount $D_i$ at times $t_i$ drops by about $D_i$ on each ex date, and no
choice of constant $q$ reproduces that. Two simple treatments are common. The \emph{escrowed} model
assumes that $S_t-\sum_{t_i>t}D_ie^{-r(t_i-t)}$, the price less the present value of the dividends still
to come, follows a geometric Brownian motion, so the call is \cref{eq5:bs} with $S$ replaced by
$S-\mathrm{PV}(D)$. The \emph{spot drop} model lets $S$ itself follow \cref{eq5:gbm} with volatility
$\sigma$ between ex dates and subtracts $D_i$ at each one. Because subtracting a fixed cash dividend
from a lognormal price can produce a negative price, this specification also needs a convention for
paths on which the stock is worth less than the dividend; here the stock is then set to zero and stays
there, which with the inputs below happened on none of the \dropPaths{} simulated paths. The two are
different models, and the spot drop version has no closed form.

Take the base option and two dividends of 2.00, at three and nine months, with present value
\divPV. The escrowed call is \escCall, against \noDivCall{} without dividends. A simulation of the spot
drop model with \dropPaths{} paths, drawn as \dropPairs{} antithetic pairs, gives \dropCall{} with
standard error \dropSE{} (computed from the pair averages), so the escrowed price is
\dropEscGap{} too low. The reason is that the escrowed model applies a 20\% volatility to the smaller
amount $S-\mathrm{PV}(D)$, so the stock's dollar volatility is too small. A common patch scales the
volatility up to $\sigma S/(S-\mathrm{PV}(D))=\sigAdj\%$; that gives \adjCall, which overshoots, and the
spot drop price less the patched price is \dropAdjGap{} at the money, \dropAdjGapLow{} at strike 70
and \dropAdjGapHigh{} at strike 130. \Cref{fig5:div} shows the gaps across strikes. The patch
overcorrects because it applies the higher volatility for the whole year, while each dividend lowers
the stock only after its ex date. Neither shortcut is exact. Across the strikes in \cref{fig5:div}
neither error exceeds \dropGapMax{} in absolute value, against an at-the-money price of \dropCall.

\begin{figure}[t]
  \centering
  \includegraphics[width=0.86\textwidth]{ch05/fig04_dividends}
  \caption{The call in the spot drop model, by simulation, less two closed form approximations, for
  strikes from 70 to 130. Two dividends of 2.00 at three and nine months; $S_0=100$, $r=5\%$,
  $\sigma=20\%$, $T=1$. The band shows two standard errors of the simulation.}
  \label{fig5:div}
\end{figure}

% =====================================================================================
\section{Checking the price by simulation}
\label{sec5:mc}

\Cref{thm5:rn} says the price is an expectation, and expectations can be estimated by averaging. For a
European call that is a check on the formula. For a payoff that depends on the path, such as an
Asian option or a fund's carried interest with a hurdle, simulation is often the most flexible
practical method, although trees, numerical integration, regression and PDE methods may also be
available.

\begin{proposition}\label{prop5:mc}
Let $Y_1,\dots,Y_n$ be independent copies of the discounted payoff $Y=e^{-rT}(S_T-K)^+$, with
$S_T=S_0\exp\bigl((r-q-\tfrac12\sigma^2)T+\sigma\sqrt TZ\bigr)$ and $Z\sim\N(0,1)$. The mean $\bar Y_n$ is
unbiased for $C$, converges to $C$ almost surely, and
$\sqrt n(\bar Y_n-C)/\widehat{\mathrm{sd}}_n\to\N(0,1)$ in distribution, where $\widehat{\mathrm{sd}}_n$
is the sample standard deviation of the $Y_i$.
Hence $\bar Y_n\pm1.96\,\widehat{\mathrm{sd}}_n/\sqrt n$ is an asymptotic 95\% confidence interval, and its width shrinks
like $n^{-1/2}$.
\end{proposition}

\begin{proof}
$\E Y=C$ by \cref{thm5:rn}. Since $0\le Y\le e^{-rT}S_T$ and $\E S_T^2<\infty$, $Y$ has a finite
variance. The strong law gives almost sure convergence, and the central limit theorem gives
$\sqrt n(\bar Y_n-C)\to\N(0,\Var Y)$. The strong law applied to $Y^2$ gives $\widehat{\mathrm{sd}}_n^2\to\Var Y$, and
Slutsky's theorem lets $\widehat{\mathrm{sd}}_n$ replace $\sqrt{\Var Y}$.
\end{proof}

For the base call the standard deviation of $Y$ is \mcSdY, so with \mcBig{} paths the standard error is
$\mcSdY/1000=\mcSEMillionTheory$. That is a slow rate: each extra digit costs a hundred times the work.
The two methods below leave the rate alone and shrink the constant in front of it.

\begin{proposition}[Antithetic variates]\label{prop5:anti}
Let $f$ be nondecreasing with $\E f(Z)^2<\infty$ and $\Var f(Z)>0$, and let $\varrho_A$ be the correlation of $f(Z)$ and
$f(-Z)$. Then $\varrho_A\le0$. Averaging $n/2$ pairs $\tfrac12(f(Z_i)+f(-Z_i))$ has variance
$(1+\varrho_A)\Var f(Z)/n$, against $\Var f(Z)/n$ for $n$ independent draws, so the efficiency gain for
the same number of payoff evaluations is $1/(1+\varrho_A)\ge1$. The standard error of the
antithetic estimate is the sample standard deviation of the $n/2$ pair averages divided by
$\sqrt{n/2}$; treating the $n$ draws as independent would overstate it.
\end{proposition}

\begin{proof}
Let $f_-(z)=f(-z)$, which is nonincreasing, and let $Z'$ be an
independent copy of $Z$. For all $z,z'$, $(f(z)-f(z'))(f_-(z)-f_-(z'))\le0$ because the two factors never
have the same strict sign. Taking expectations, $2\E[f(Z)f_-(Z)]-2\E f(Z)\E f_-(Z)\le0$, which is
$\operatorname{Cov}(f(Z),f(-Z))\le0$. Since $-Z$ has the law of $Z$, both terms of a pair have
variance $\Var f(Z)$, so the pair average has variance $\tfrac12\Var f(Z)(1+\varrho_A)$; dividing by the
$n/2$ pairs gives the result.
\end{proof}

\begin{proposition}[Control variates]\label{prop5:cv}
Let $X$ be a random variable with known mean and finite positive variance, here $X=e^{-rT}S_T$ with $\E_QX=S_0e^{-qT}$. For a
constant $\beta$ the variable $Y_\beta=Y-\beta(X-\E X)$ has mean $C$, and its variance is smallest at
$\beta^*=\operatorname{Cov}(X,Y)/\Var X$, where it equals $(1-\varrho^2)\Var Y$ with $\varrho$ the correlation of
$X$ and $Y$.
\end{proposition}

\begin{proof}
$\Var Y_\beta=\Var Y-2\beta\operatorname{Cov}(X,Y)+\beta^2\Var X$ is a quadratic in $\beta$ minimised at $\beta^*$, and
substituting gives $\Var Y-\operatorname{Cov}(X,Y)^2/\Var X=(1-\varrho^2)\Var Y$.
\end{proof}

In practice $\beta^*$ is estimated from the same paths, by the sample covariance over the sample
variance. That introduces a bias of order $1/n$ but leaves the asymptotic variance unchanged
\parencite[Section 4.1]{glasserman2003}. The standard error is then the sample standard deviation of
the residuals $Y_i-\hat\beta(X_i-\E X)$, with the estimated $\hat\beta$, divided by $\sqrt n$. Because
$\hat\beta$ is estimated, this is a plug-in standard error, valid as $n$ grows; it is not exact in
finite samples. For the
base call, exact moments of the lognormal give $\varrho=\mcRho$ for the control and
$\varrho_A=\mcCorrAnti$ for the antithetic pair, so the predicted gains are \vrCVTheory{} for the control and
\vrAntiTheory{} for the antithetic pair.

With \mcBig{} payoff evaluations each (for the antithetic estimators, \mcPairs{} pairs), the four
estimators give
\begin{center}
\begin{tabular}{lcccc}
  \toprule
  & plain & antithetic & control variate & both\\
  \midrule
  estimate & \mcEstPlain & \mcEstAnti & \mcEstCV & \mcEstBoth\\
  standard error & \mcSEPlain & \mcSEAnti & \mcSECV & \mcSEBoth\\
  variance reduction & 1 & \vrMCAnti & \vrMCCV & \vrMCBoth\\
  \bottomrule
\end{tabular}
\end{center}
against the exact \mcExact. The measured gains match the predicted ones. Used together the two methods
reduce the variance by a factor of \vrMCBoth, more than \vrMCProduct, the product of the separate gains, and
close to the \vrBothTheory{} predicted from exact moments. The antithetic average has already removed
the part of the payoff that is odd in $Z$. What is left, the even part, has correlation
\mcRhoEven{} with the even part of the control, against \mcRho{} for the raw payoff and the raw
control.
Over \mcReps{} independent repetitions at each of five sample sizes from 250 to 64,000 payoff
evaluations, the root mean
square error falls with slopes \mcSlopePlain, \mcSlopeAnti, \mcSlopeCV{} and \mcSlopeBoth{} on a
log-log scale, as \cref{prop5:mc} predicts, and the 95\% intervals at $n=4{,}000$ contain the true price
in \mcCoverPlain, \mcCoverAnti, \mcCoverCV{} and \mcCoverBoth{} percent of repetitions
(\cref{fig5:mc}).

\begin{figure}[t]
  \centering
  \includegraphics[width=\textwidth]{ch05/fig05_mc}
  \caption{(a) A plain Monte Carlo estimate of the base call as paths are added, with a band of two
  standard errors. (b) Root mean square error over \mcReps{} repetitions against the number of payoff
  evaluations for the four estimators, with the $n^{-1/2}$ reference line.}
  \label{fig5:mc}
\end{figure}

The listing below is the whole of a variance reduced pricer, short enough to retype.
\lstinputlisting[language=Python]{code/ch05/snippets/bs_mc.py}
It uses seed \snipSeed{} and \snipPairs{} antithetic pairs with the stock as a control, and prints a
formula value of \snipFormula{} and a simulated \snipMC{} with standard error \snipSE. The simulated
price is \snipZ{} standard errors below the formula, computed from the unrounded output. For a normal estimator a miss at least that large
in either direction has probability \snipProb\%, about one run in \snipOdds. The seed was not changed
after the first run, so the result is reported as it came out.

\paragraph{Greeks by simulation}
Sensitivities can be estimated from the same paths. For the call, two estimators are unbiased for
delta.

\begin{proposition}\label{prop5:pathwise}
With $S_T$ as in \cref{prop5:mc},
\begin{equation*}
  \Delta=\E_Q\Bigl[e^{-rT}\mathbf 1\{S_T>K\}\frac{S_T}{S_0}\Bigr]
  =\E_Q\Bigl[e^{-rT}(S_T-K)^+\frac{Z}{S_0\sigma\sqrt T}\Bigr].
\end{equation*}
The first is the \emph{pathwise} estimator, the second the \emph{likelihood ratio} estimator.
\end{proposition}

\begin{proof}
For the pathwise form, $S_T$ is proportional to $S_0$, so $\partial S_T/\partial S_0=S_T/S_0$, and
$s\mapsto(s-K)^+$ is Lipschitz with derivative $\mathbf 1\{s>K\}$ except at $s=K$, which has
probability zero. The difference quotients $\bigl((S_T(S_0+\delta S)-K)^+-(S_T(S_0)-K)^+\bigr)/\delta S$ are bounded
by $S_T/S_0$, which is integrable, so dominated convergence lets the derivative pass inside the
expectation. For the likelihood ratio form, write the price as $\int e^{-rT}(s-K)^+p(s;S_0)\,ds$, where
$p$ is the lognormal density of $S_T$. With $z(s)=\bigl(\ln(s/S_0)-(r-q-\tfrac12\sigma^2)T\bigr)/(\sigma\sqrt T)$,
$\partial p/\partial S_0=p(s;S_0)\,z(s)/(S_0\sigma\sqrt T)$. To differentiate under the integral it is
enough that this derivative, times $(s-K)^+$, is bounded by an integrable function of $s$ uniformly
for $S_0$ in a compact interval $[S_0',S_0'']\subset(0,\infty)$. On that interval
$|z(s)|\le c_2(1+|\ln s|)$ and $p(s;S_0)\le c_3\,s^{-1}\exp\bigl(-(\ln s)^2/(4\sigma^2T)\bigr)$ for all
large $|\ln s|$, with constants $c_2,c_3$ that depend only on the interval, and
$(s-K)^+(1+|\ln s|)s^{-1}\exp\bigl(-(\ln s)^2/(4\sigma^2T)\bigr)$ is integrable on $(0,\infty)$. So
dominated convergence applies and gives the second form.
\end{proof}

With \deltaPaths{} paths the pathwise estimate is \deltaPW{} with standard error \deltaPWSE, and the
likelihood ratio estimate \deltaLR{} with standard error \deltaLRSE, against the exact \deltaC. The
pathwise standard error is \deltaSERatio{} times smaller here, about \deltaVarRatio{} times less variance.
The likelihood ratio estimator is the one to use when the payoff jumps, as a digital option's does.
There the pathwise derivative is zero almost everywhere, so the pathwise estimator returns zero while
the true delta is positive, because all of the sensitivity sits at the jump.

% =====================================================================================
\section{Checking the price on a grid}
\label{sec5:fd}

The third way to get a price is to solve \cref{eq5:pde} on a grid. In log price $x=\ln S$ and time to
expiry $\tau$, the equation is $V_\tau=aV_{xx}+bV_x-rV$ with constant coefficients $a=\tfrac12\sigma^2$
and $b=r-q-\tfrac12\sigma^2$. The scripts use a uniform grid of $M$ intervals centred on $\ln S_0$ and
reaching six standard deviations, $6\sigma\sqrt T$, on each side, with the values $0$ and
$Se^{-q\tau}-Ke^{-r\tau}$ for a call imposed at the two ends. Write $V^k_j$ for the value at
$x_j=x_0+j\,\Delta x$ and $\tau_k=k\,\Delta\tau$, and $L_h$ for the central difference operator
\begin{equation*}
  (L_hV)_j=a\frac{V_{j+1}-2V_j+V_{j-1}}{\Delta x^2}+b\frac{V_{j+1}-V_{j-1}}{2\Delta x}-rV_j .
\end{equation*}
For a weight $w\in[0,1]$ the \emph{weighted scheme} is
\begin{equation}\label{eq5:theta}
  \frac{V^{k+1}-V^k}{\Delta\tau}=w\,L_hV^{k+1}+(1-w)L_hV^k,
\end{equation}
which is \emph{explicit} for $w=0$, \emph{implicit} for $w=1$ and \emph{Crank-Nicolson} for
$w=\tfrac12$ \parencite{crank1947}. Many texts call it the $\theta$-scheme; here $\theta$ already names the market price of risk.
For $w>0$ each step solves a tridiagonal system, which costs a number of operations proportional to
$M$.

\begin{lemma}\label{lem5:heat}
Let $\gamma_1=-b/(2a)$ and $\gamma_2=-b^2/(4a)-r$. Then $V(\tau,x)=e^{\gamma_1x+\gamma_2\tau}u(\tau,x)$ solves
$V_\tau=aV_{xx}+bV_x-rV$ if and only if $u_\tau=au_{xx}$.
\end{lemma}

\begin{proof}
With $\mathcal E=e^{\gamma_1x+\gamma_2\tau}$, $V_\tau=\mathcal E(\gamma_2u+u_\tau)$, $V_x=\mathcal E(\gamma_1u+u_x)$ and
$V_{xx}=\mathcal E(\gamma_1^2u+2\gamma_1u_x+u_{xx})$. Substituting and dividing by $\mathcal E$,
$u_\tau=au_{xx}+(2a\gamma_1+b)u_x+(a\gamma_1^2+b\gamma_1-r-\gamma_2)u$. The first bracket is zero by the choice of
$\gamma_1$, and the second is $b^2/(4a)-b^2/(2a)-r-\gamma_2=0$ by the choice of $\gamma_2$.
\end{proof}

So the stability of the scheme can be studied on the heat equation, where the drift and discounting
are gone. Apply \cref{eq5:theta} to $u_\tau=au_{xx}$ and set $\lambda=a\,\Delta\tau/\Delta x^2$.

\begin{theorem}[Von Neumann stability]\label{thm5:vonneumann}
The Fourier mode $u^k_j=A^ke^{\mathrm ij\omega}$ solves the weighted scheme for the heat equation if and
only if
\begin{equation}\label{eq5:amp}
  A(\omega)=\frac{1-4(1-w)\lambda\chi}{1+4w\lambda\chi},\qquad \chi=\sin^2(\omega/2).
\end{equation}
Every mode stays bounded, $|A(\omega)|\le1$ for all $\omega$, if and only if $w\ge\tfrac12$, or
$w<\tfrac12$ and $\lambda\le1/\bigl(2(1-2w)\bigr)$. In particular the explicit scheme needs
$\lambda\le\tfrac12$, the implicit and Crank-Nicolson schemes are stable for every $\lambda$, and as
$\lambda\to\infty$ the highest frequency $\omega=\pi$ has $A\to0$ for the implicit scheme and $A\to-1$ for
Crank-Nicolson.
\end{theorem}

\begin{proof}
For the mode, $u_{j+1}-2u_j+u_{j-1}=(e^{\mathrm i\omega}-2+e^{-\mathrm i\omega})u_j=-4\chi\,u_j$. Substituting
in \cref{eq5:theta} gives $A-1=-4\lambda\chi\bigl(wA+(1-w)\bigr)$, which rearranges to
\cref{eq5:amp}. Since $\lambda\chi\ge0$, the numerator never exceeds the denominator, so $A\le1$. The
condition $A\ge-1$ is $1-4(1-w)\lambda\chi\ge-1-4w\lambda\chi$, that is $2\lambda\chi(1-2w)\le1$.
If $w\ge\tfrac12$ this holds for all $\lambda$. If $w<\tfrac12$ the worst case is $\chi=1$, at
$\omega=\pi$, which gives the bound on $\lambda$. At $\chi=1$, $A=(1-4(1-w)\lambda)/(1+4w\lambda)$,
which tends to $-(1-w)/w$ as $\lambda\to\infty$: zero for $w=1$ and $-1$ for $w=\tfrac12$.
\end{proof}

The theorem is exact for the heat equation on an infinite grid. For the explicit scheme applied to
$V$ directly, the same substitution gives
\begin{equation*}
  A(\omega)=1-4\lambda\chi-r\Delta\tau+\mathrm i\,\frac{b\Delta\tau}{\Delta x}\sin\omega ,\qquad
  \Bigl(\frac{b\Delta\tau}{\Delta x}\Bigr)^2=\frac{b^2\lambda\,\Delta\tau}{a}.
\end{equation*}
If $\lambda\le\tfrac12$ then $|1-4\lambda\chi|\le1$, so the real part has absolute value at most
$1+|r|\Delta\tau$, for a rate of either sign, and $|A|^2\le(1+|r|\Delta\tau)^2+b^2\lambda\Delta\tau/a$. For
$\Delta\tau\le1$ this gives $|A|\le1+c_1\Delta\tau$ with $c_1=|r|+b^2\lambda/(2a)+r^2/2$, the drift Courant
number $b\Delta\tau/\Delta x$ and the discount rate $r$ entering only through $c_1$. This does not show
$|A|\le1$ for the explicit scheme with drift; it bounds growth over a fixed horizon, which is the form
of stability that matters, because $(1+c_1\Delta\tau)^{T/\Delta\tau}\le e^{c_1T}$. For $w\ge\tfrac12$ and $r\ge0$ the factor is
$A=\bigl(1-(1-w)z\bigr)/(1+wz)$ with $z=4\lambda\chi+r\Delta\tau-\mathrm i(b\Delta\tau/\Delta x)\sin\omega$.
The real part of the denominator is at least the absolute value of the real part of the numerator, and the
imaginary part of the denominator is at least as large in absolute value as that of the numerator, so
$|A|\le1$ with no restriction on $\lambda$.

\begin{proposition}[Consistency]\label{prop5:consistency}
For a solution with bounded derivatives of the orders that appear, the local truncation error of
\cref{eq5:theta} is $O(\Delta\tau+\Delta x^2)$, and $O(\Delta\tau^2+\Delta x^2)$ when $w=\tfrac12$.
\end{proposition}

\begin{proof}
By Taylor's theorem,
\begin{equation*}
  \frac{V_{j+1}-V_{j-1}}{2\Delta x}=V_x+O(\Delta x^2),\qquad
  \frac{V_{j+1}-2V_j+V_{j-1}}{\Delta x^2}=V_{xx}+O(\Delta x^2),
\end{equation*}
so $L_hV=LV+O(\Delta x^2)$, where $L$ is the
differential operator. Expanding about $\tau^*=\tau_k+w\Delta\tau$, the left side of
\cref{eq5:theta} is $V_\tau(\tau^*)+(\tfrac12-w)\Delta\tau\,V_{\tau\tau}+O(\Delta\tau^2)$ and
$w\,LV^{k+1}+(1-w)LV^k=LV(\tau^*)+O(\Delta\tau^2)$, because $\tau^*$ is the matching weighted
average of $\tau_k$ and $\tau_{k+1}$. Since $V_\tau=LV$, the error is
$(\tfrac12-w)\Delta\tau\,V_{\tau\tau}+O(\Delta\tau^2+\Delta x^2)$, and the first term vanishes
exactly when $w=\tfrac12$.
\end{proof}

The Lax equivalence theorem of \textcite{lax1956}, which we quote, says that a consistent scheme
for a well posed linear problem converges if and only if it is stable. It says nothing about the
rate: an order of convergence needs a separate error bound, and that bound needs the solution to be
smooth enough. Here it is not. The call payoff has a kink, so the derivatives in
\cref{prop5:consistency} are not bounded at expiry, and the grid is truncated at six standard
deviations with boundary values that are exact only far from the strike. We prove no order, and the
rates below are measurements. They are measured in one setting: the strike on a grid node, the grid
centred on $S_0$ and truncated at $6\sigma\sqrt T$, and the error taken at $S_0$. In that setting the
explicit and Crank-Nicolson prices converge at second order in $M$ and the implicit price at first
order with $M/2$ steps, and \cref{fig5:fdstab} shows that the same is not true of every Greek at
every node.

\Cref{fig5:fdconv} shows the error at $S_0$ as $M$ grows, with $\lambda=0.4$ fixed for the explicit
scheme and $M/2$ time steps for the others. The fitted orders in $M$ are \fdOrderExpl{} for explicit,
\fdOrderImpl{} for implicit and \fdOrderCN{} for Crank-Nicolson. The implicit scheme is first order because
$\Delta\tau$ shrinks only as fast as $\Delta x$; the explicit scheme is second order in $M$ because
$\lambda$ fixed forces $\Delta\tau\propto\Delta x^2$. At $M=400$ the errors are \fdErrExpl{} with
\fdStepsExpl{} steps for explicit, \fdErrImpl{} for implicit and \fdErrCN{} for Crank-Nicolson with
\fdStepsCN{} steps; Crank-Nicolson gives \fdCNValue{} against \callBase. The explicit scheme looks best,
but it takes \fdStepRatio{} times as many steps, and it is partly lucky. Crank-Nicolson's error barely moves when its
steps are raised to \fdStepsCNFine, where it is \fdErrCNFine, so it is almost all spatial. The explicit
scheme carries the same spatial error plus a time error of the opposite sign, and at $\lambda=0.4$ the
two largely cancel.

\begin{figure}[t]
  \centering
  \includegraphics[width=0.86\textwidth]{ch05/fig06_fd_convergence}
  \caption{Error of the finite difference price of the base call at $S_0$ against the number of space
  intervals, on a log-log scale. Explicit at $\lambda=0.4$; the other schemes, including
  Crank-Nicolson with a Rannacher start, with $M/2$ time steps. Grid centred on $\ln S_0$ and
  truncated at $6\sigma\sqrt T$, with the strike on a node.}
  \label{fig5:fdconv}
\end{figure}

The stability limit is sharp: just inside it the explicit scheme is accurate, and just outside it fails at
once. Each step count below is rounded up, so the Courant number never exceeds the value named. With
$M=200$ it takes, at $\lambda=0.45$, \explStepsA{}
steps and is off by \explErrA. At $\lambda=0.50$, the edge of the region, it takes \explStepsB{} steps and
the error rises to \explErrB. Asking for $\lambda=0.52$ gives only \explStepsC{} steps and an actual $\lambda=\explCourC$; the highest frequency mode
grows by a factor of $|1-4\lambda|=\explAmpC$ a step and the price is off by \explErrC.

Crank-Nicolson's weakness is the other end of \cref{thm5:vonneumann}. With large steps, $A(\pi)$ is
close to $-1$, so the high frequency content of the kinked payoff flips sign each step and is barely
damped. The price at $S_0$ hardly notices, but gamma does. For a three month option on a grid of
\cnSpace{} intervals with \cnSteps{} time steps, $\lambda$ is about \cnLambda, and Crank-Nicolson's gamma
for stock prices from 80 to 125, the range plotted in \cref{fig5:fdstab}(b), is off by up to
\gammaErrCN, against a true peak of \gammaPeakExact. The cure, due to \textcite{rannacher1984} in the
form recommended by \textcite{giles2006}, is to replace the first two steps with four half steps of
the implicit scheme, whose $A(\pi)$ tends to zero. That reduces the gamma error to \gammaErrCNR{} while keeping second
order in time afterwards; the implicit scheme alone has an error of \gammaErrImpl{} (\cref{fig5:fdstab}).

\begin{figure}[t]
  \centering
  \includegraphics[width=\textwidth]{ch05/fig07_fd_stability}
  \caption{(a) Von Neumann amplification factors \cref{eq5:amp}. The explicit scheme leaves $[-1,1]$
  for $\lambda>\tfrac12$; Crank-Nicolson tends to $-1$ at high frequencies when $\lambda$ is large.
  (b) Gamma of a three month call from Crank-Nicolson with \cnSteps{} large steps on \cnSpace{} space
  intervals, with and without a Rannacher start, against the exact gamma.}
  \label{fig5:fdstab}
\end{figure}

% =====================================================================================
\section{Reading volatility from option prices}
\label{sec5:iv}

Spot, strike, maturity and the interest rate in \cref{eq5:bs} are observed, while volatility has to
be inferred from option prices; dividend forecasts and funding rates may need estimating too. Markets
therefore quote options by the volatility that reproduces their price.

\begin{definition}\label{def5:iv}
The \emph{implied volatility} of a European option with market price $C^{\mathrm{mkt}}$ is the $\sigma>0$ at which the
Black-Scholes price equals $C^{\mathrm{mkt}}$.
\end{definition}

\begin{proposition}\label{prop5:ivexist}
A call price $C^{\mathrm{mkt}}$ has an implied volatility if and only if
$(Se^{-q\tau}-Ke^{-r\tau})^+<C^{\mathrm{mkt}}<Se^{-q\tau}$, and it is then unique. A put with the same strike and
expiry, priced consistently with \cref{eq5:parity}, has the same implied volatility.
\end{proposition}

\begin{proof}
By \cref{prop5:parity} the call price is continuous and strictly increasing in $\sigma$ and runs from
the lower bound to the upper bound as $\sigma$ goes from $0$ to $\infty$; the intermediate value
theorem gives existence and monotonicity gives uniqueness. The two sides of \cref{eq5:parity} do not
depend on $\sigma$, so a call and a put that satisfy it are matched by the same $\sigma$.
\end{proof}

There is no closed form for the inverse, so it is found numerically. Writing the price in forward
units makes the problem one of two variables. Let $y=\ln(F/K)$ and $\nu=\sigma\sqrt\tau$ and define
\begin{equation}\label{eq5:normcall}
  c(y,\nu)=\frac{C}{e^{-r\tau}F}=\Phi\Bigl(\frac{y}{\nu}+\frac{\nu}{2}\Bigr)-e^{-y}\,\Phi\Bigl(\frac{y}{\nu}-\frac{\nu}{2}\Bigr),
\end{equation}
with $d_1=y/\nu+\nu/2$ and $d_2=y/\nu-\nu/2$ as before.

\begin{lemma}\label{lem5:shape}
$\partial c/\partial \nu=\varphi(d_1)$ and $\partial^2c/\partial \nu^2=\varphi(d_1)d_1d_2/\nu$. So for $y\ne0$,
$c$ is strictly convex in $\nu$ on $(0,\nu^*)$ and strictly concave on $(\nu^*,\infty)$, where
$\nu^*=\sqrt{2|y|}$; for $y=0$ it is strictly concave on $(0,\infty)$. The normalised put
$c-(1-e^{-y})$ has the same derivatives in $\nu$.
\end{lemma}

\begin{proof}
In these units \cref{lem5:phi} reads $\varphi(d_1)=e^{-y}\varphi(d_2)$. Then
\begin{equation*}
  \frac{\partial c}{\partial \nu}=\varphi(d_1)\frac{\partial d_1}{\partial \nu}-e^{-y}\varphi(d_2)\frac{\partial d_2}{\partial \nu}
  =\varphi(d_1)\frac{\partial(d_1-d_2)}{\partial \nu}=\varphi(d_1),
\end{equation*}
and differentiating again,
\begin{equation*}
  \frac{\partial^2c}{\partial \nu^2}=-d_1\varphi(d_1)\frac{\partial d_1}{\partial \nu}=-d_1\varphi(d_1)\Bigl(\frac12-\frac{y}{\nu^2}\Bigr)
  =\frac{\varphi(d_1)d_1d_2}{\nu},
\end{equation*}
using $d_2/\nu=y/\nu^2-\tfrac12$. Finally
$d_1d_2=y^2/\nu^2-\nu^2/4$, which is positive for $\nu<\nu^*$ and negative for $\nu>\nu^*$. The put differs from
the call by a term that does not involve $\nu$.
\end{proof}

\begin{theorem}[Monotone convergence of Newton's method away from forward at-the-money]\label{thm5:newton}
Let $y\ne0$, and let $\hat\nu>0$ be the unique solution of $c(y,\hat\nu)=\hat c$, with $\hat c$ strictly between
the no-arbitrage price bounds. Set $\nu^*=\sqrt{2|y|}>0$, and apply Newton's method
$\nu_{k+1}=\nu_k-[c(y,\nu_k)-\hat c]/\varphi(d_1(\nu_k))$, starting from $\nu_0=\nu^*$. If $\hat\nu\ge\nu^*$, the
iterates increase monotonically to $\hat\nu$; if $\hat\nu<\nu^*$, they decrease monotonically to $\hat\nu$. In
either case, convergence is eventually quadratic. When $y=0$, the inflection-point start is $\nu^*=0$, where
the displayed expression for $d_1$ is undefined; this starting rule is therefore not covered by the
theorem.
\end{theorem}

\begin{proof}
Let $f(\nu)=c(y,\nu)-\hat c$, which is strictly increasing. Suppose $\hat\nu\ge \nu^*$ and $\nu^*\le \nu_k\le\hat\nu$,
so $f(\nu_k)\le0$ and $\nu_{k+1}\ge \nu_k$. On $[\nu^*,\infty)$, $f$ is concave, so its tangent at $\nu_k$ lies
above its graph; the tangent vanishes at $\nu_{k+1}$, so $f(\nu_{k+1})\le0$ and $\nu_{k+1}\le\hat\nu$. By
induction the sequence increases and is bounded by $\hat\nu$, so it has a limit $\bar\nu$, and passing to
the limit in the iteration gives $f(\bar\nu)=0$, hence $\bar\nu=\hat\nu$. If $\hat\nu<\nu^*$ the argument is the
mirror image on $(0,\nu^*]$, where $f$ is convex and the tangent lies below the graph, so the iterates
decrease and stay above $\hat\nu>0$. The tangent argument needs each iterate to stay in $(0,\nu^*]$. It
does: $f$ extends continuously to $\nu=0$ with $f(0)=(1-e^{-y})^+-\hat c<0$, because $\hat c$ exceeds the
intrinsic value. By convexity on $[0,\nu^*]$ the tangent at $\nu_k$ is negative at $0$, so it vanishes at
a positive point. Quadratic convergence near $\hat\nu$ is the standard property of Newton's method for a
twice continuously differentiable function with $f'(\hat\nu)\ne0$.
\end{proof}

The starting point $\nu^*$ is the inflection point of the price, a choice due to \textcite{manaster1982}.
At the forward, where $y=0$, the code starts at zero and uses the limits there, $c=0$ and slope
$\varphi(0)$; the theorem says nothing about that start, and a separate test checks it.
A start elsewhere can fail. Where the price curve is nearly flat, at very small or very large $\nu$,
the tangent is nearly horizontal and the first step can shoot to a huge volatility or a negative one. The code works with the
out-of-the-money option, using parity to convert a quote if needed. By \cref{lem5:shape} the iterates
are the same either way, but a deep in-the-money price is mostly intrinsic value, and subtracting that
from the target loses digits that the out-of-the-money price keeps. For production work,
\textcite{jaeckel2015} gives an inversion built from a carefully constructed first approximation and
a fixed number of high-order refinement steps, designed to reach close to machine precision.

Take a six month call struck at 120 with $\sigma=30\%$, price \ivExPrice. The start is
$\sigma_0=\nu^*/\sqrt\tau=\ivExStart$, above the answer, so the iterates fall: \ivExStepA, \ivExStepB,
\ivExStepC. The iteration stops when a step changes $\nu$ by less than $10^{-12}$ (relative to $\nu$
once $\nu>1$). The fourth iterate is already within \ivExErrFour{} of 0.3, and the fifth step, smaller
than $10^{-12}$, triggers the stop, so the count is \ivExIter{} iterations. Brent's method \parencite{brent1973}
on $[0.0001,5]$, which needs no derivative and always keeps a bracket, takes \ivExBrentIter{} iterations
and \ivExBrentCalls{} price evaluations (\cref{fig5:newton}). Over a grid of \ivGridN{} cases, calls and
puts at three expiries, four volatilities from 5\% to 100\% and strikes within one standard deviation of
the forward, both methods recover the volatility to within \ivGridWorst. Newton needs at most
\ivNewtonMax{} iterations (median \ivNewtonMed) and its iterates are monotone in every case; Brent needs
at most \ivBrentMax{} (median \ivBrentMed).

\begin{figure}[t]
  \centering
  \includegraphics[width=\textwidth]{ch05/fig08_newton}
  \caption{(a) The normalised call price against total volatility for $\ln(F/K)=-0.5$, with Newton steps
  from the inflection point $\nu^*$ to two targets, one on each side. (b) Error against the number of
  evaluations for the 120 strike example; each Newton step costs a price and a vega.}
  \label{fig5:newton}
\end{figure}

\begin{proposition}[At the money]\label{prop5:bsapprox}
At the money forward, $y=0$ and $c(0,\nu)=2\Phi(\nu/2)-1=\dfrac{\nu}{\sqrt{2\pi}}\Bigl(1-\dfrac{\nu^2}{24}\Bigr)+O(\nu^5)$.
So $\sqrt{2\pi}\,C/(e^{-r\tau}F\sqrt\tau)$ approximates $\sigma$ with relative error $-\nu^2/24+O(\nu^4)$.
\end{proposition}

\begin{proof}
At $y=0$, $d_1=\nu/2=-d_2$, and \cref{eq5:normcall} gives $\Phi(\nu/2)-\Phi(-\nu/2)=2\int_0^{\nu/2}\varphi$. Expand
$\varphi(z)=\varphi(0)(1-z^2/2+O(z^4))$ and integrate: $2\varphi(0)\bigl(\nu/2-\nu^3/48\bigr)+O(\nu^5)$, and
$2\varphi(0)=\sqrt{2/\pi}$.
\end{proof}

This is the rule of thumb of \textcite{brenner1988}, and it is also the first step of the iteration the code uses at the forward, started at
zero.
For a one year option struck at the forward, \bsApproxStrike, with $\sigma=20\%$, it gives \bsApprox\%.
Applied to the base call, which is at the money spot and so in the money forward, the same formula
gives \bsApproxSpot\%: the rule needs the strike at the forward.

\paragraph{Why one volatility does not fit every strike}
If the model were right, every option on a stock would have the same implied volatility. They do not.
Equity index options show a skew, with low strikes at higher volatilities, and the pattern is
steeper at short maturities. Jump diffusion and stochastic volatility can both generate smiles and
skews, although neither reproduces every feature of an observed equity index surface.

\begin{proposition}[Jump diffusion]\label{prop5:merton}
Let the stock follow \cref{eq5:gbm} under $Q$ with drift $r-q-\lambda_Jk_J$, plus jumps at the times of an
independent Poisson process $N^J$ of rate $\lambda_J$ that multiply the price by $e^{J_i}$ with
$J_i\sim\N(\mu_J,\delta_J^2)$ independent, and $k_J=e^{\mu_J+\delta_J^2/2}-1$. Then
\begin{equation*}
\begin{gathered}
  C=\sum_{i=0}^\infty e^{-\lambda_J'T}\frac{(\lambda_J'T)^i}{i!}\,C_{\mathrm{BS}}(\sigma_i,r_i),\\
  \lambda_J'=\lambda_J(1+k_J),\qquad\sigma_i^2=\sigma^2+\frac{i\delta_J^2}T,\qquad r_i=r-\lambda_Jk_J+\frac{i\ln(1+k_J)}T,
\end{gathered}
\end{equation*}
where $C_{\mathrm{BS}}(\sigma_i,r_i)$ is \cref{eq5:bs} with those inputs \parencite{merton1976}.
\end{proposition}

\begin{proof}
The drift correction makes $\E_QS_T=S_0e^{(r-q)T}$, since $\E e^{J_i}=1+k_J$ and the expected number of
jumps is $\lambda_JT$. Condition on $N^J_T=i$. Then
$\ln S_T\sim\N\bigl(\ln S_0+(r-q-\lambda_Jk_J-\tfrac12\sigma^2)T+i\mu_J,\ \sigma^2T+i\delta_J^2\bigr)$. A
lognormal with rate $r_i$ and volatility $\sigma_i$ has the same variance and the mean
$\ln S_0+(r_i-q-\tfrac12\sigma_i^2)T$, which equals the one above because $\ln(1+k_J)=\mu_J+\tfrac12\delta_J^2$.
So $e^{-rT}\E_Q[(S_T-K)^+\mid N^J_T=i]=e^{(r_i-r)T}C_{\mathrm{BS}}(\sigma_i,r_i)$. Multiplying by
$Q(N^J_T=i)=e^{-\lambda_JT}(\lambda_JT)^i/i!$ and using $e^{(r_i-r)T}=e^{-\lambda_Jk_JT}(1+k_J)^i$ gives the
weights shown.
\end{proof}

The proposition specifies the jump intensity and the jump-size distribution under $Q$. How they relate
to the jump parameters under $\mathbb P$ is a separate modelling assumption. \textcite{merton1976} took them
to be the same by assuming that jump risk is diversifiable and so carries no premium, and nothing in
the market model forces that choice: with jumps, the stock and the bond can no longer replicate every
claim, the hedge argument of \cref{thm5:pde} breaks down, and many measures $Q$ are consistent with
the absence of arbitrage.

The second extension makes the variance itself random. In the model of \textcite{heston1993},
\begin{equation*}
  dS_t=(r-q)S_t\,dt+\sqrt{v_t}\,S_t\,dW^1_t,\qquad dv_t=\kappa(\bar v-v_t)\,dt+\xi\sqrt{v_t}\,dW^2_t,\qquad
  d\langle W^1,W^2\rangle_t=\varrho\,dt .
\end{equation*}
A negative $\varrho$ makes volatility rise when the price falls, which fattens the left tail. The
characteristic function of $\ln S_T$ is known in closed form, and we quote the single integral of
\textcite{lewis2000}:
\begin{equation*}
  C=e^{-rT}\Bigl(F-\frac{\sqrt{FK}}\pi\int_0^\infty\operatorname{Re}\bigl[e^{\mathrm i\omega\ln(F/K)}\psi(\omega-\tfrac{\mathrm i}2)\bigr]\frac{d\omega}{\omega^2+\frac14}\Bigr),
\end{equation*}
where $F=S_0e^{(r-q)T}$, $T$ is the time to expiry in years, and
$\psi(u)=\E_Q\bigl[e^{\mathrm iu\ln(S_T/F)}\bigr]$ is the characteristic function of $\ln(S_T/F)$ under $Q$,
extended to $u=\omega-\tfrac{\mathrm i}2$, where it is finite because $\E_QS_T<\infty$. It is coded in the
form of \textcite{albrecher2007} that avoids jumps in the complex logarithm, and the integral is
computed by adaptive quadrature on $(0,\infty)$ (SciPy's \texttt{quad} with absolute tolerance $10^{-14}$,
relative tolerance $10^{-12}$ and at most 1,000 subintervals). With a Gaussian $\psi$ the same code returns
\cref{eq5:bs} to $10^{-12}$, which is one of the tests. A full truncation Euler simulation
\parencite{lord2010} with \hesMCPaths{} paths and \hesMCSteps{} steps gives \hesMC{} (standard error
\hesMCSE) for a six month at-the-money-forward call for which the integral gives \hesFourier. The gap is
\hesMCZ{} standard errors, so at this sample size the scheme's time discretisation bias is too small
to detect.

\Cref{fig5:smiles} shows the two smiles. The jump model uses $\sigma=15\%$, $\lambda_J=0.5$ jumps a year,
$\mu_J=-10\%$ and $\delta_J=15\%$; the stochastic volatility model uses $v_0=\bar v=0.04$, $\kappa=2$,
$\xi=0.5$ and $\varrho=-0.7$. Measured by the volatility at $\ln(K/F)=-0.1$ less that at $+0.1$, the jump
model's skew is \skewMerShort, \skewMerMid{} and \skewMerLong{} points at expiries of 0.1, 0.5 and 2
years, and the stochastic volatility model's is \skewHesShort, \skewHesMid{} and \skewHesLong. Jumps give a
steep short dated skew that flattens fast, because over a long horizon many small jumps look like a
diffusion; the jump model's at-the-money volatility rises from \atmMerShort\% to \atmMerLong\% toward
the volatility of the log return including jumps, $\sqrt{\sigma^2+\lambda_J(\mu_J^2+\delta_J^2)}=\merTotalVol\%$. Stochastic volatility
keeps its skew longer. Its at-the-money volatility, \atmHesShort\%, \atmHesMid\% and \atmHesLong\%, sits
below the 20\% that $v_0=\bar v=0.04$ suggests, because at the money the price is concave in variance
and averaging over random variance lowers it, an argument that is exact for $\varrho=0$
\parencite{hullwhite1987} and approximate here. \Cref{fig5:surface} shows the stochastic volatility surface,
which runs from \hesSurfMin\% to \hesSurfMax\%.

\begin{figure}[t]
  \centering
  \includegraphics[width=\textwidth]{ch05/fig09_smiles}
  \caption{Implied volatility against log moneyness at three expiries. (a) Merton jump diffusion.
  (b) Heston stochastic volatility.}
  \label{fig5:smiles}
\end{figure}

\begin{figure}[t]
  \centering
  \includegraphics[width=0.8\textwidth]{ch05/fig10_heston_surface}
  \caption{The implied volatility surface of the stochastic volatility model, for log moneyness from
  $-0.4$ to $0.3$ and expiries from 0.1 to 2 years.}
  \label{fig5:surface}
\end{figure}

\paragraph{What a smile says about the distribution}
Arbitrage-free European call prices at one expiry, known for a continuum of strikes, determine the
risk-neutral distribution of $S_T$, whatever model produced them.

\begin{proposition}[Breeden-Litzenberger]\label{prop5:bl}
Let $C(K)=e^{-rT}\E_Q(S_T-K)^+$, and suppose $S_T$ has a continuous density $p$ under $Q$. Then
$\partial C/\partial K=-e^{-rT}Q(S_T>K)$ and $\partial^2C/\partial K^2=e^{-rT}p(K)$
\parencite{breeden1978}.
\end{proposition}

\begin{proof}
$C(K)=e^{-rT}\int_K^\infty(s-K)p(s)\,ds$. Differentiating in $K$, the boundary term vanishes because the
integrand is zero at $s=K$, leaving $-e^{-rT}\int_K^\infty p(s)\,ds$; differentiating once more gives
$e^{-rT}p(K)$ by the fundamental theorem of calculus.
\end{proof}

\begin{corollary}\label{cor5:static}
Call prices in a market without arbitrage are decreasing and convex in the strike, with slope in
$[-e^{-rT},0]$. A payoff $g$ that is twice continuously differentiable can be written, with $F$ the
forward price (the expansion holds about any positive point, and the forward is the usual choice), as
\begin{equation*}
  g(S_T)=g(F)+g'(F)(S_T-F)+\int_0^Fg''(K)(K-S_T)^+\,dK+\int_F^\infty g''(K)(S_T-K)^+\,dK ,
\end{equation*}
so a claim on $g(S_T)$ is a bond, a forward and a strip of puts and calls \parencite{carr1998}.
\end{corollary}

\begin{proof}
The first part reads the signs off \cref{prop5:bl}, since $0\le Q(S_T>K)\le1$ and $p\ge0$. For the
second, Taylor's theorem with integral remainder gives, for $s>F$,
$g(s)=g(F)+g'(F)(s-F)+\int_F^sg''(K)(s-K)\,dK$, and the last integral equals
$\int_F^\infty g''(K)(s-K)^+dK$, while the put integral is zero; the case $s<F$ is symmetric.
\end{proof}

Applied to the stochastic volatility prices at six months, with strikes from 20 to 260 in steps of 0.1,
the second difference of the call prices integrates to \densIntRaw, the discount factor $e^{-rT}=\densDisc$.
Multiplied by $e^{rT}$, as \cref{prop5:bl} requires, it is the density under $Q$: it integrates to
\densInt{} and has mean equal to the forward,
\densFwd, to within \densTol. \Cref{fig5:density} compares the density with a lognormal at the same at-the-money
volatility, \densAtmVol\%. The model puts \tailHes\% of probability below 80\% of the forward against
\tailLogn\% for the lognormal, and \rtailHes\% above 120\% against \rtailLogn\%. Under $Q$ the skew gives
a crash more weight and a rally less than the lognormal does; prices alone cannot say how much of that
is real-world probability and how much is a premium for crash risk.

\begin{figure}[t]
  \centering
  \includegraphics[width=0.86\textwidth]{ch05/fig11_density}
  \caption{The risk-neutral density of $S_T$ implied by the stochastic volatility prices at six months,
  by \cref{prop5:bl}, against a lognormal density with the same at-the-money volatility.}
  \label{fig5:density}
\end{figure}

% =====================================================================================
\section{What a hedger actually earns}
\label{sec5:pnl}

\Cref{thm5:pde} assumes the hedger knows $\sigma$ and trades continuously. Neither is true. This
section works out what a hedger earns when the volatility is wrong and when the hedge is adjusted only
at discrete times.

\begin{theorem}[Hedging at the wrong volatility]\label{thm5:pnl}
Let the stock follow $dS_t=\mu_tS_t\,dt+\sigma_tS_t\,dW_t$, $S_0>0$, and pay the constant yield $q$,
where $\mu_t$ and $\sigma_t$ are progressively measurable, unknown to the hedger, and satisfy
$\int_0^T(|\mu_t|+\sigma_t^2)\,dt<\infty$ almost surely, so that $S$ is a positive continuous It\^o
process and It\^o's lemma applies to $V(t,S_t)$ on $[0,T)$. For a volatility $\sigma$ let $V(\cdot;\sigma)$ be the
solution of \cref{eq5:pde} with that volatility and payoff $g$, with the properties required in
\cref{thm5:pde}. A dealer sells the claim for $V^p_0=V(0,S_0;\sigma_p)$ and holds
$\Delta_t=V_S(t,S_t;\sigma_h)$ shares, self-financing. Write $\Gamma^h_t=V_{SS}(t,S_t;\sigma_h)$ and
assume
\begin{equation}\label{eq5:pnlint}
  \int_0^T\Gamma^h_tS_t^2\,\bigl|\sigma_h^2-\sigma_t^2\bigr|\,dt<\infty\quad\text{almost surely.}
\end{equation}
Then the discounted profit at $T$, after the hedge is liquidated at $S_T$ and the claim paid, is
\begin{equation}\label{eq5:pnl}
  \mathrm{P\&L}=V(0,S_0;\sigma_p)-V(0,S_0;\sigma_h)+\int_0^Te^{-rt}\,\tfrac12\Gamma^h_tS_t^2\bigl(\sigma_h^2-\sigma_t^2\bigr)dt .
\end{equation}
\end{theorem}

\begin{proof}
Let $\Pi$ be the value of the dealer's hedge portfolio, with $\Pi_0=V^p_0$, and write $V^h_t=V(t,S_t;\sigma_h)$.
As in the proof of \cref{thm5:rn}(ii), $d(e^{-rt}\Pi_t)=e^{-rt}\Delta_t\bigl(dS_t+(q-r)S_t\,dt\bigr)$. By
It\^o's lemma with the true volatility and the equation for $V^h$,
\begin{equation*}
\begin{aligned}
  d(e^{-rt}V^h_t)&=e^{-rt}\bigl(V_t-rV+\tfrac12\sigma_t^2S_t^2V_{SS}\bigr)dt+e^{-rt}V_S\,dS_t\\
  &=e^{-rt}\bigl(\tfrac12(\sigma_t^2-\sigma_h^2)S_t^2\Gamma^h_t-(r-q)S_t\Delta_t\bigr)dt+e^{-rt}\Delta_t\,dS_t ,
\end{aligned}
\end{equation*}
using $V_t-rV=-\tfrac12\sigma_h^2S^2V_{SS}-(r-q)SV_S$. Subtracting,
$d\bigl(e^{-rt}(\Pi_t-V^h_t)\bigr)=e^{-rt}\tfrac12\Gamma^h_tS_t^2(\sigma_h^2-\sigma_t^2)\,dt$. Integrate from
$0$ to $t<T$ and let $t\to T$. The left side converges because $\Pi$ and $V^h$ are continuous up to
$T$, with $V^h_T=g(S_T)$, and the right side converges by \cref{eq5:pnlint}. The profit is
$e^{-rT}(\Pi_T-g(S_T))$, which is \cref{eq5:pnl}. The identity holds path by path, so it needs no
admissibility condition on the hedge; whether the dealer could be forced out before $T$ by a lower
bound on wealth is a separate question.
\end{proof}

Condition \cref{eq5:pnlint} is where the behaviour of the volatility enters. For a call or a put
$\Gamma^h_tS_t^2=S_te^{-q\tau}\varphi(d_1)/(\sigma_h\sqrt\tau)\le\max(1,e^{-qT})\,S_t/(\sigma_h\sqrt{2\pi\tau})$,
since $\tau\le T$ and $\varphi\le1/\sqrt{2\pi}$; the factor is one when $q\ge0$. Then
$\tau^{-1/2}$ is integrable on $[0,T]$, so, since each path of $S$ is continuous and hence bounded on
$[0,T]$, the condition holds whenever $\sigma_t$ is bounded.

\begin{corollary}[Robustness]\label{cor5:robust}
If the dealer prices and hedges at the same $\sigma_h=\sigma_p$, the profit is
$\int_0^Te^{-rt}\tfrac12\Gamma^h_tS_t^2(\sigma_h^2-\sigma_t^2)\,dt$. If, in addition, the payoff $g$ is
convex and $\sigma_t\le\sigma_h$ at all times, the profit is nonnegative on every
path \parencite{elkaroui1998}.
\end{corollary}

\begin{proof}
The first terms of \cref{eq5:pnl} cancel. A convex payoff has a convex Black-Scholes price, because
$V(t,s)=e^{-r\tau}\E_Qg(sG)$ with $G=S_T/S_t$ lognormal and independent of $S_t$, and
$s\mapsto g(sG)$ is convex for each outcome of $G$; so $\Gamma^h\ge0$ and the integrand is
nonnegative.
\end{proof}

So a dealer who prices and hedges a convex claim at the same volatility earns a nonnegative discounted
profit on every path if the instantaneous volatility never exceeds that hedge volatility, as long as
the hedge is continuous; a realised volatility that is lower only on average is not enough. How much depends on the path, through the dollar gamma
$\Gamma S^2$: a large move in volatility counts for little if it happens when the option is far from the
money.

\begin{proposition}\label{prop5:qmean}
Suppose the true volatility is a constant $\sigma$ and the drift is $\mu$. Then
\begin{enumerate}[label=(\roman*)]
  \item if $\sigma_h=\sigma$, the profit is exactly $V(0,S_0;\sigma_p)-V(0,S_0;\sigma)$ on every path;
  \item for any $\sigma_h$, the expected profit under $Q$ is $V(0,S_0;\sigma_p)-V(0,S_0;\sigma)$.
\end{enumerate}
\end{proposition}

\begin{proof}
(i) The integral in \cref{eq5:pnl} is zero. (ii) The discounted hedge portfolio is a $Q$-martingale by
\cref{thm5:rn}(ii), since $\Delta$ is bounded and so $\E_Q\int_0^T\Delta_t^2S_t^2\,dt<\infty$; hence
$\E_Q[e^{-rT}\Pi_T]=V^p_0$, and
$\E_Q[e^{-rT}g(S_T)]=V(0,S_0;\sigma)$ by \cref{thm5:bs}.
\end{proof}

The proposition separates two choices. Hedging at the true volatility locks in the mispricing; hedging
at the implied volatility leaves the same expected profit under $Q$ but makes it depend on the path.
\textcite{ahmad2005} compare the two choices in detail.
Under the real-world measure, with $\mu\ne r$, the expected profit of the second choice is no longer
pinned down, because the path's dollar gamma depends on where the stock drifts.

The simulations in \texttt{hedging\_pnl.py} use \hedgePaths{} paths, the base inputs and daily
rebalancing unless stated otherwise. A call is sold at 20\% on a stock
whose volatility turns out to be 15\%, and the gap in value is \volGap. Hedged daily at 15\%, the
profit has mean \wrongRealMean{} and standard deviation \wrongRealSd, the leftover coming from discrete
rebalancing. Hedged at 20\%, the mean is \wrongImpMean{} (standard error \wrongImpSE) and the standard
deviation \wrongImpSd. Hedged at 20\%, \wrongImpLoss\% of paths lose money, as \cref{cor5:robust} leads one
to expect, though that no path loses is an outcome of this finite simulation and nothing guarantees
it path by path once the hedge is rebalanced discretely; hedged at 15\%, \wrongRealLoss\% do, through the rebalancing error. With a real-world drift of \hedgeMu\% instead
of $r$, hedging at 20\% gives a mean of \wrongImpMeanDrift: the stock tends to drift away from the strike,
where gamma is smaller. A dealer who sells at 20\% when 25\% is realised, and hedges at 20\%, loses
\underLoss{} on average (standard error \underSE), against $V(25\%)-V(20\%)=\underGapAbs$
(\cref{fig5:gt}(b) shows the 20\% and 15\% case).

\paragraph{Hedging at intervals}
Over one rebalancing interval $\Delta t$, the dealer's short call, hedged at the right volatility, gains
$e^{-rt}\tfrac12\Gamma S^2(\sigma^2\Delta t-R^2)$ to leading order, where $R=S_{t+\Delta t}/S_t-1$ is the
stock's simple return over the interval; the approximation drops drift, financing and higher-order
terms in $R$, all of smaller order in $\Delta t$. This is the gamma-theta relation of \cref{prop5:gammatheta} in discrete time: for the
short option, theta pays $\tfrac12\Gamma S^2\sigma^2\Delta t$ and gamma costs $\tfrac12\Gamma S^2R^2$,
and they cancel only on average. \Cref{fig5:gt}(a) plots the daily profit on one path against the stock's return; the points lie
on the parabola, and their correlation with the gamma-theta term is \gtStepCorr. Over the year the
profit totals \gtTotal, against \gtApprox{} from the sum of the gamma-theta terms.

\begin{proposition}[Discrete hedging error]\label{prop5:discrete}
With the true volatility $\sigma$, hedging a call at $N$ equally spaced dates $\Delta t=T/N$ leaves a
discounted profit with mean exactly zero under $Q$ and, to leading order in $\Delta t$, variance
\begin{equation}\label{eq5:hedgevar}
  \frac{\Delta t}2\int_0^Te^{-2rt}\sigma^4\,\E\bigl[(\Gamma_tS_t^2)^2\bigr]dt ,
\end{equation}
where the expectation and the variance are taken under the same measure, $Q$ or $\mathbb P$. The
measure matters for the variance through the distribution of the stock path; under $\mathbb P$ the mean
is not zero in general.
\end{proposition}

\begin{proof}[Derivation]
The statement about the mean is exact. The discounted value of any self-financing portfolio, here one with
a piecewise constant and bounded $\Delta$, is a $Q$-martingale by \cref{thm5:rn}(ii), and the discounted
payoff has $Q$-mean $V(0,S_0;\sigma)$, which was the premium. The variance is a leading-order
calculation. Over a step, $R=\sigma\sqrt{\Delta t}\,Z+O(\Delta t)$ with $Z$ standard normal, so
$\sigma^2\Delta t-R^2=\sigma^2\Delta t(1-Z^2)+O(\Delta t^{3/2})$. Given the past, the step profit has
conditional mean of order $\Delta t^{3/2}$ and conditional variance
$\tfrac14e^{-2rt}\Gamma^2S^4\sigma^4\Delta t^2\Var(Z^2)=\tfrac12e^{-2rt}\Gamma^2S^4\sigma^4\Delta t^2$, since
$\Var(Z^2)=2$. The leading terms of different steps are uncorrelated because each has conditional mean
zero, so their variances add; the cancellation of cross-step covariances is shown only for these
leading martingale-difference terms, and the remainders are of higher order. The sum over $N$ steps is
a Riemann sum for \cref{eq5:hedgevar}. The
integral is finite because, by \cref{lem5:gammamoment}, $\E[(\Gamma_tS_t^2)^2]$ grows only like
$(T-t)^{-1/2}$ near expiry. \textcite{boyle1980} is an early study of discrete hedging, and a rigorous
version of the expansion is given by \textcite{bertsimas2000}.
\end{proof}

\begin{lemma}\label{lem5:gammamoment}
For a call or a put, if $\ln S_t\sim\N(m,\eta^2)$, then with $\varsigma^2=\tfrac12\sigma^2\tau$ and
$\ell=\ln K-(r-q+\tfrac12\sigma^2)\tau$,
\begin{equation*}
  \E\bigl[(\Gamma_tS_t^2)^2\bigr]=\frac{e^{-2q\tau}\varsigma}{\sqrt{2\pi}\,\sigma^2\tau}\,e^{2m+2\eta^2}\,
  \frac{1}{\sqrt{\eta^2+\varsigma^2}}\,\varphi\Bigl(\frac{m+2\eta^2-\ell}{\sqrt{\eta^2+\varsigma^2}}\Bigr).
\end{equation*}
\end{lemma}

\begin{proof}
By \cref{prop5:greeks}, $\Gamma S^2=Se^{-q\tau}\varphi(d_1)/(\sigma\sqrt\tau)$ with $d_1=(\ln S-\ell)/(\sigma\sqrt\tau)$.
So $(\Gamma S^2)^2=e^{-2q\tau}S^2e^{-d_1^2}/(2\pi\sigma^2\tau)$, and
$e^{-d_1^2}=e^{-(\ln S-\ell)^2/(2\varsigma^2)}=\sqrt{2\pi}\,\varsigma\,n_{\varsigma^2}(\ln S-\ell)$, where
$n_a(x)=\varphi(x/\sqrt a)/\sqrt a$ is the normal density with variance $a$ (used only in this proof).
For a value $x$ of $\ln S_t$, $e^{2x}$ times the $\N(m,\eta^2)$ density equals $e^{2m+2\eta^2}$ times the
$\N(m+2\eta^2,\eta^2)$ density, and the convolution of normal densities gives
\begin{equation*}
  \int n_{\eta^2}(x-m-2\eta^2)\,n_{\varsigma^2}(x-\ell)\,dx=n_{\eta^2+\varsigma^2}(m+2\eta^2-\ell).\qedhere
\end{equation*}
\end{proof}

The standard deviation of the hedging error therefore falls like $N^{-1/2}$: to halve it, rebalance four
times as often. For the base call with a real-world drift of \hedgeMu\%, the simulated and leading-order
standard deviations are
\begin{center}
\begin{tabular}{lcccc}
  \toprule
  rebalances a year & \hedgeNThirteen & \hedgeNFiftyTwo & \hedgeNDaily & \hedgeNQuarter\\
  \midrule
  simulated & \hedgeSdThirteen & \hedgeSdFiftyTwo & \hedgeSdDaily & \hedgeSdQuarter\\
  leading order \cref{eq5:hedgevar} & \hedgeThThirteen & \hedgeThFiftyTwo & \hedgeThDaily & \hedgeThQuarter\\
  simulated, percent of the premium & \hedgeSdPctThirteen & \hedgeSdPctFiftyTwo & \hedgeSdPctDaily & \hedgeSdPctQuarter\\
  \bottomrule
\end{tabular}
\end{center}
The fitted slope on a log-log scale is \hedgeSlope. The leading-order formula overstates the error a
little at low frequencies, where the higher-order terms matter, and is within \hedgeThGapDaily\% at daily
hedging. The mean profits are all within \hedgeMeanMaxZ{} standard errors of zero. The paths use the
real-world drift, so this is a finding for the chosen drift and not a consequence of
\cref{prop5:discrete}, which guarantees a zero mean only under $Q$. The correlation of the profit with the
summed gamma-theta terms is \hedgeCorrGTFiftyTwo{} for weekly and \hedgeCorrGTDaily{} for daily hedging.
\Cref{fig5:hedge} shows the distributions.

\begin{figure}[t]
  \centering
  \includegraphics[width=\textwidth]{ch05/fig12_gamma_theta}
  \caption{(a) Daily profit of a hedged short call along one path against the stock's return, with the
  gamma-theta term, with daily rebalancing. (b) Sold at 20\% when 15\% is realised: profit when hedged
  daily at 20\% and at 15\%, \hedgePaths{} paths, drift $r$.}
  \label{fig5:gt}
\end{figure}

\begin{figure}[t]
  \centering
  \includegraphics[width=0.86\textwidth]{ch05/fig13_hedge_hist}
  \caption{Discounted profit of a short call hedged at the true volatility, for \hedgeNThirteen,
  \hedgeNFiftyTwo{} and \hedgeNDaily{} rebalances a year. \hedgePaths{} paths, base inputs, real-world
  drift \hedgeMu\%.}
  \label{fig5:hedge}
\end{figure}


\paragraph{What rebalancing costs}
If each trade costs a fraction $\varepsilon$ of its value, the cost of rebalancing grows with the frequency. The
change in delta over a step is about $\Gamma S\sigma\sqrt{\Delta t}\,|Z|$, so a step costs about
$\varepsilon\Gamma S^2\sigma\sqrt{\Delta t}\,\E|Z|=\varepsilon\Gamma S^2\sigma\sqrt{\Delta t}\sqrt{2/\pi}$, and the total over
$N$ steps grows like $\sqrt N$. With $\varepsilon=\costRate\%$, the expected rebalancing cost is \costRebalThirteen,
\costRebalFiftyTwo{} and \costRebalDaily{} at 13, 52 and 252 rebalances a year, a slope of \costSlope{} on a
log-log scale, on top of \costInit{} for the initial hedge. More frequent hedging reduces the leading-order variance of hedge error but increases expected
transaction costs; the actual costs still vary with the price path.

\textcite{leland1985} suggested charging for the cost through the volatility. The expected cost per
unit time, $\varepsilon\Gamma S^2\sigma\sqrt{2/\pi}/\sqrt{\Delta t}$, has the same form as the gain
$\tfrac12\Gamma S^2(\sigma_L^2-\sigma^2)$ from hedging at a higher volatility $\sigma_L$, and equating the
two gives
\begin{equation*}
  \sigma_L^2=\sigma^2\Bigl(1+\sqrt{\frac2\pi}\,\frac{2\varepsilon}{\sigma\sqrt{\Delta t}}\Bigr).
\end{equation*}
For weekly hedging $\sigma_L=\lelandVol\%$ and the call costs \lelandPrice, a markup of \lelandMarkup.
Sold and hedged at $\sigma_L$ with costs, the mean profit is \lelandMean{} (standard error \lelandSE,
standard deviation \lelandSd). That is almost exactly the cost of the opening trade, which the argument
ignores; without it the mean is \lelandExInit. Under this proportional cost model the markup
approximately funds the rebalancing. It is not a no-arbitrage price once costs are present, since
with costs the hedge no longer replicates the option exactly.

% =====================================================================================
\section{Case study: selling US dollar calls in Canada}
\label{sec5:case}

The study takes the dealer's side. At each month end it sells a one month call on one US dollar, struck
at the forward, and hedges it daily at the Bank of Canada's published rate until the next month end.
It is a historical, frictionless hedging experiment on published benchmark rates, not a strategy a
dealer could have executed as written: there are no bid-offer spreads, and the published rates are
not dealing prices. The
volatility charged and used for the hedge is the realised volatility of the previous \fxWindow{} trading days,
so the dealer has no information about the future. The domestic rate is the Canadian one month
Treasury bill yield and the foreign rate the US one month bill yield, each converted from a simple
30 day yield on an actual/365 basis to a continuous rate. The option is priced and hedged with
\cref{prop5:black}(ii). A month end is the last day in the calendar month with a published rate. The
option's life is the number of calendar days to the next month end divided by 365, and the hedge is
adjusted on each day with a published rate. Both volatilities are computed from daily log changes,
annualised by the calendar time they span. Profits are discounted to the sale date and expressed in
basis points of spot.

The data file has \fxDays{} days with a published exchange rate, from \fxStart{} to \fxEnd. The file is cut at the end of
September 2026, and a month enters the study only once it is over. After
the first \fxWindow{} days that gives \fxN{} options sold from \fxFirst{} to \fxLast; the last is
closed on \fxLastClose, the final rate published that month, and every option runs for a full
month. The premium averaged \fxPremMean{} basis points.

The seller earned a mean of \fxPnlMean{} basis points a month, with standard error \fxPnlSE, median
\fxPnlMedian{} and standard deviation \fxPnlSd. On average the price was fair, within the noise. The
hedging volatility averaged \fxShMean\% against \fxSrMean\% realised over the option's life, and it was
above the realised volatility in \fxShAbove\% of months. The seller made money in \fxWin\% of months.

The gamma-theta terms of \cref{thm5:pnl} track most of the month to month variation. The correlation of each month's profit with its
summed gamma-theta terms is \fxCorrGT. Within this backtest, the close relationship is consistent with the gamma-theta mechanism: volatility
differences explain much of the profit variation, while discrete rebalancing and the timing of price
moves leave an additional residual. A correlation alone does not prove the explanation. A regression of the profit on
$\mathcal V\times(\sigma_h-\sigma_r)$, where $\sigma_h$ is the hedging volatility and $\sigma_r$ the
volatility realised over the month, measures the first-order effect of the volatility gap. It has slope
\fxSlope{} (standard error \fxSlopeSE) and a coefficient of determination of \fxRsq. The slope is below one partly because the vega approximation
ignores the path: a volatile day counts for more when spot is near the strike, which the gamma-theta sum
records and the monthly volatility does not.

The worst month was the option sold at the end of \fxWorstMonth, which lost \fxWorstLoss{} basis points when
volatility rose from \fxWorstSh\% over the previous three months to \fxWorstSr\% over the next month. The best
was \fxBest{} basis points, for the option sold at the end of \fxBestMonth{} when the hedging volatility
still carried the turmoil of that autumn. The five worst months account for \fxTopFiveShare\% of all
losses, and without them the mean rises to \fxPnlMeanExWorst. Over the whole period the profits sum to
\fxCum{} basis points of spot. Selling volatility at its trailing value roughly broke even over these \fxN{} months, with
losses that came in clusters when volatility jumped (\cref{fig5:case}).

\begin{figure}[t]
  \centering
  \includegraphics[width=\textwidth]{ch05/fig14_case}
  \caption{(a) USD/CAD volatility over the previous \fxWindow{} trading days, used to price and hedge each
  option sold from \fxFirst{} to \fxLast, and the volatility realised over its life. (b) Each month's hedged profit against the sum of its
  gamma-theta terms, in basis points of spot, with the 45 degree line.}
  \label{fig5:case}
\end{figure}

One caution about the data. The Bank published a noon rate, a snapshot, until 28 April 2017. Since
1 March 2017 it has published a single daily average \parencite{bocfxchange}, and its Valet series
carries that average back to 3 January 2017; the study switches from the noon series to the average
at the start of 2017, so the option sold at the end of December 2016 is hedged on both series, and the
trailing volatilities for options sold from January to March 2017 mix them. An average of prices is smoother than the prices, and differences
of averages are positively autocorrelated even when the underlying changes are not \parencite{working1960}.
The lag one autocorrelation of daily log changes is \acRetNoon{} in the noon series and \acRetAvg{} in
the average series. Since both the hedging and the realised volatility are computed from the same
series, the comparison is internally consistent, but a dealer could not have traded at an average
rate, so the second half of the study is further from an executable hedge than the first. The
mean profit is similar in both periods: \fxPnlMeanNoon{} basis points over \fxNNoon{} months on noon
rates and \fxPnlMeanAvg{} over \fxNAvg{} months on average rates.

% =====================================================================================
\section{What the model leaves out}
\label{sec5:limits}

The derivation used every assumption in \cref{def5:market}. Real markets violate each of them, sometimes
in ways that matter for pricing or hedging.

Volatility is not constant. The smile in \cref{sec5:iv} is the market's way of saying so, and in the
case study the realised volatility went from \fxWorstSh\% over one quarter to \fxWorstSr\% over the
next month. A dealer who hedges at
the wrong volatility carries the profit in \cref{eq5:pnl}, which is a bet on realised volatility weighted
by dollar gamma.

Prices jump. With jumps, the stock and the bank account cannot replicate an option, there is no unique
price without an assumption about how jump risk is priced, and a delta hedge leaves a residual that
does not shrink as rebalancing becomes more frequent.

Trading is discrete and costs money. \Cref{prop5:discrete} and the cost estimates give the trade-off,
and in practice dealers rebalance when the hedge drifts outside a band, whatever the date. Borrowing costs more than lending,
and shares cannot always be borrowed for a short hedge, so a single rate $r$ gives way to separate
borrowing and lending rates, and one price gives way to a band between the costs of hedging a long and
a short position.

Rates and dividends are uncertain. How much that matters depends on maturity, moneyness, the
underlying's sensitivity to rates and the hedging instruments available. For a short dated equity
option near the money it is usually a small effect beside volatility; for a ten year option or a
convertible bond, rate risk is no longer small beside equity risk.

For alternative assets the deepest problem is that the underlying may not trade at all. A real option
on a development site, a manager's carry or an option on a private company's equity has no hedge in the
underlying, so the argument of \cref{thm5:pde} does not apply. The formula can still be used as a
valuation by analogy, but its inputs then carry the burden: the volatility is an estimate of a
quantity no one observes, and the risk-neutral drift is an assumption that the market price of risk
of the asset is known. The hedging error results of \cref{sec5:pnl} assume a traded underlying, so
for these assets they describe the size of the problem but do not measure it.

% =====================================================================================
\section{Summary}
\label{sec5:summary}

A delta hedge removes the stock's random term and leaves an equation without the expected return. The
same price is a discounted expectation under the measure in which the stock earns the riskless rate
less its yield. Solving gives the formula, whose Greeks follow from one identity between the two
normal densities, and whose theta and gamma are tied together by the equation itself. Futures,
currencies and indices fit by choosing $q$; cash dividends need an approximation. Monte Carlo
converges like $n^{-1/2}$, and antithetic and control variates cut its variance by a factor of
\vrMCBoth{} here. Crank-Nicolson is unconditionally stable and second order in our measurements but needs a damped start for
smooth Greeks, and the explicit scheme needs $\lambda\le\tfrac12$. Newton's method from the inflection
point inverts the formula monotonically. A hedger who uses the wrong volatility earns the gap between
hedging and realised variance, weighted by dollar gamma; a hedger who rebalances $N$ times a year
carries an error of order $N^{-1/2}$; and
\fxYears{} years of US dollar options written in Canada show the first effect in data, with monthly
profits that track the summed gamma-theta terms.

% =====================================================================================
\section{Exercises}
\label{sec5:exercises}

\begin{enumerate}[label=\textbf{5.\arabic*}, leftmargin=*, ref=5.\arabic*]
  \item\label{ex5:one} A stock trades at 50, pays no dividend and has volatility 25\%; the rate is 4\%.
  Price six month calls and puts struck at 55, and check put-call parity.
  \item\label{ex5:two} For the call in Exercise~\ref{ex5:one}, estimate the price after the stock rises by 1 using
  delta alone and using delta and gamma, and compare with the exact price.
  \item\label{ex5:three} At what stock price is the gamma of the option in Exercise~\ref{ex5:one} largest?
  \item\label{ex5:four} A futures contract trades at 80. Price a three month at-the-money call with
  $r=3\%$ and $\sigma=30\%$, and compare with the approximation of \cref{prop5:bsapprox}.
  \item\label{ex5:five} Spot is 1.36 Canadian dollars per US dollar, the Canadian rate 3\% and the US rate
  4.5\%, both continuous, and volatility 8\%. Price a six month at-the-money spot put on one million US
  dollars, and give the forward rate.
  \item\label{ex5:six} Show that a digital call paying 1 if $S_T>K$ is worth $e^{-r\tau}\Phi(d_2)$, and
  that this equals $-\partial C/\partial K$.
  \item\label{ex5:seven} The explicit scheme is run in log price with $\Delta x=0.01$ and $\sigma=30\%$.
  What is the largest stable time step, and how many steps does a one year option need?
  \item\label{ex5:eight} A nine month call struck at 105 on a stock at 100 trades at 8.00 with $r=3\%$
  and no dividends. Find its implied volatility with Newton's method from $\nu^*$.
  \item\label{ex5:nine} For the base call hedged weekly, with the \hedgeMu\% real-world drift of the
  table in \cref{sec5:pnl}, \cref{eq5:hedgevar} gives a standard deviation of \xvNineSd. How many rebalances a year bring it down to 0.25?
  \item\label{ex5:ten} A dealer sells the base call at 25\% when the volatility turns out to be 20\%. What
  is the profit if the dealer hedges continuously at 20\%? What can be said if the hedge uses 25\%?
\end{enumerate}

\subsection*{Answers}

\begin{enumerate}[label=\textbf{5.\arabic*}, leftmargin=*]
  \item $d_1=\xvOneDOne$ and $d_2=\xvOneDTwo$; the call is worth \xvOneCall{} and the put \xvOnePut.
  $C-P=\xvOneParity=50-55e^{-0.02}$, as \cref{eq5:parity} requires.
  \item $\Delta=\xvTwoDelta$ and $\Gamma=\xvTwoGamma$. Delta alone gives \xvTwoDeltaOnly, delta and gamma
  \xvTwoApprox, and the exact price is \xvTwoExact.
  \item By \cref{prop5:gammapeak}, $S^*=55e^{-(0.04+1.5\times0.25^2)\times0.5}=\xvThreeS$.
  \item By \cref{prop5:black}(i) the call is \xvFourCall. The approximation
  $e^{-rT}F\sigma\sqrt T/\sqrt{2\pi}$ gives \xvFourApprox, slightly high, as the factor $1-\nu^2/24$ predicts.
  \item The forward is $1.36e^{(0.03-0.045)\times0.5}=\xvFiveFwd$, and \cref{prop5:black}(ii) gives a put
  worth \xvFivePut{} Canadian dollars.
  \item The digital is worth $e^{-r\tau}Q(S_T>K)=e^{-r\tau}\Phi(d_2)$ by \cref{lem5:lognormal}.
  Differentiating \cref{eq5:bs} in $K$ as in the proof of \cref{prop5:greeks}, the terms in
  $\partial d_i/\partial K$ cancel by \cref{lem5:phi}, leaving $\partial C/\partial K=-e^{-r\tau}\Phi(d_2)$.
  This is also the first part of \cref{prop5:bl}.
  \item $\lambda=\tfrac12\sigma^2\Delta\tau/\Delta x^2\le\tfrac12$ gives $\Delta\tau\le\Delta x^2/\sigma^2=\xvSevenDt$,
  so \xvSevenSteps{} steps for one year.
  \item $y=\ln(100e^{0.0225}/105)$, $\nu^*=\sqrt{2|y|}$, so the start is $\sigma_0=\xvEightStart\%$. The
  first step gives \xvEightFirst\%, and the iteration stops after \xvEightIter{} steps at \xvEightSig\%.
  The start is already close because the option is nearly at the money forward.
  \item The standard deviation scales as $N^{-1/2}$, so $N=52\times(\xvNineSd/0.25)^2$, rounded up to
  \xvNineN, about \xvNinePerDay{} a trading day. The formula at that frequency gives \xvNineCheck.
  \item By \cref{prop5:qmean}(i), hedging at the true 20\% locks in $V(25\%)-V(20\%)=\xvTenGap$ on every
  path. Hedging at 25\%, the profit depends on the path, but by \cref{cor5:robust} it is nonnegative on
  every path and by \cref{prop5:qmean}(ii) its expectation under $Q$ is the same \xvTenGap.
\end{enumerate}

% =====================================================================================
\section{Notes, sources and credits}
\label{sec5:credits}

\paragraph{Sources}
The model and the hedging argument are from \textcite{black1973} and \textcite{merton1973}, and the
martingale approach from \textcite{harrison1981}; \textcite{shreve2004} develops both in the style used
here. The futures formula is from \textcite{black1976} and the currency formula from
\textcite{garman1983}. Monte Carlo pricing of options goes back to \textcite{boyle1977}; \textcite{glasserman2003} is the standard reference for variance reduction and Greeks.
The finite difference schemes go back to \textcite{crank1947}, the damped start to
\textcite{rannacher1984}, in the form analysed by \textcite{giles2006}, and the equivalence theorem to
\textcite{lax1956}; \textcite{wilmott2006} gives a
practical treatment. The Newton start is from \textcite{manaster1982}, the at-the-money approximation from
\textcite{brenner1988} and a fast exact inversion from \textcite{jaeckel2015}; Brent's method is from
\textcite{brent1973}. The jump model is from \textcite{merton1976}, the stochastic volatility model from
\textcite{heston1993}, its single integral from \textcite{lewis2000} and the stable characteristic
function from \textcite{albrecher2007}, and the full truncation scheme from \textcite{lord2010}; the
concavity argument for the at-the-money level is from \textcite{hullwhite1987}, and
\textcite{gatheral2006} surveys the volatility surface. The density result is from
\textcite{breeden1978} and the spanning formula from \textcite{carr1998}. The robustness result is from \textcite{elkaroui1998}, the comparison of hedging volatilities follows
\textcite{ahmad2005}, the
discrete hedging error from \textcite{boyle1980} and \textcite{bertsimas2000}, and the cost adjusted volatility from
\textcite{leland1985}. The averaging effect is from \textcite{working1960}. \textcite{hull2022} covers market
conventions.

\paragraph{Software used}
The scripts call NumPy \parencite{numpy} for arrays and random numbers, SciPy \parencite{scipy} for the normal
distribution, banded solvers, quadrature and Brent's method, pandas \parencite{pandas} for the case study
data, statsmodels \parencite{statsmodels} for its regression and Matplotlib \parencite{matplotlib} for the figures. The R
cross-check uses base R \parencite{rcore} and the Octave cross-check uses the quadrature, sparse
solvers and root finders of GNU Octave \parencite{octave}. The references give the version of each
package used for the numbers printed here, dated by the release year of that version, and
\texttt{code/requirements-lock.txt} records the full environment; a test checks that the installed
versions match the references.

\paragraph{Projects for comparison}
The code here does not call these, but they implement the same methods in production form and are
useful for checking: QuantLib \parencite{quantlib} has closed forms, finite difference engines and the
stochastic volatility model, and py\_vollib \parencite{pyvollib} has fast implied volatility inversion. The full
list, with links, is in \texttt{CREDITS.md}.

\paragraph{Data}
The case study uses four daily series. From the Bank of Canada Valet service \parencite{bocfx}: the
noon rate for the US dollar, series \texttt{IEXE0101}, up to 30 December 2016; the daily average
rate, \texttt{FXUSDCAD}, from 3 January 2017; and the one month Treasury bill yield,
\texttt{TB.CDN.30D.MID}. From FRED \parencite{fredDGS}: the US one month constant maturity Treasury
yield, \texttt{DGS1MO}. Rates are Canadian dollars per US dollar and yields are in percent. The series were
retrieved on 19 October 2026 and cut at 30 September 2026, so the file covers \fxStart{} to \fxEnd,
the last closing date of the study, as explained in \cref{sec5:case}. Days without a published exchange rate are dropped; a
missing yield on a day with a rate is replaced by the previous day's. The two exchange rate series are
joined without adjustment at the turn of 2017, and the yields are converted to continuous rates as
described in \cref{sec5:case}.
The script \texttt{scripts/build\_usdcad.py} rebuilds \texttt{data/usdcad\_rates.csv} from the raw
downloads in \texttt{data/raw/}, cut at the same date, and a test checks that the rebuilt file is identical. The Bank can
revise published values, and the numbers here use the saved files.

\printbibliography[heading=subbibliography,title={References for Chapter \thechapter}]
\end{refsection}
